% ThermalQualityRevision20261023-v1
% QualityRevision20261023
\hypersetup{bookmarksdepth=3}
\begingroup
\renewcommand{\Needspace}[1]{\par\begingroup
\dimen0=#1\relax\vskip0pt plus\dimen0\penalty-100
\vskip0pt plus-\dimen0\vskip\dimen0\penalty9999\vskip-\dimen0
\vskip0pt\endgroup}
\tcbset{breakable/.default=false}
\chapter{流体静力学与动力学}
% Editable TikZ: control volumes, free bodies, and local integration elements.
\providecommand{\ThermalFigure}[3]{\begin{qualityfigure}
\begin{tikzpicture}[>=Latex,font=\small,x=.8cm,y=.8cm]#1\end{tikzpicture}
\figcap{#2}\label{#3}\end{qualityfigure}}
\newcommand{\ThFlowElement}[1]{
 \draw[thick](0,.5)--(2.5,.5)--(4,1)--(6.5,1);
 \draw[thick](0,-.5)--(2.5,-.5)--(4,-1)--(6.5,-1);
 \fill[DeepBlue!18](.6,-.5)rectangle(1.8,.5);\fill[DeepBlue!18](4.5,-1)rectangle(5.1,1);
 \draw[<->](.6,-1.5)--node[below]{$v_A\Delta t$}(1.8,-1.5);
 \draw[<->](4.5,-1.5)--node[below]{$v_B\Delta t$}(5.1,-1.5);
 \node[above]at(1.2,.7){$S_A$};\node[above]at(4.8,1.2){$S_B$};#1}
\newcommand{\ThAFlowVolume}{\ThermalFigure{\ThFlowElement{}}{在同一时间内经过两个截面的流体柱，长度分别为当地流速乘时间。}{ch13:quality:volume-a}
图~\ref{ch13:quality:volume-a}中的阴影柱体不是同一时刻被跟踪的两个质点，而是相同时间内穿过截面的体积。稳态、不可压缩且无侧向支流保证控制体没有净体积积累，因而两柱体积相等。若右侧截面变大，右柱必须变短，正对应流速减小。
}
\newcommand{\ThAUniformFlow}{\ThermalFigure{
 \draw[thick](0,-.5)rectangle(7,.5);\draw[dashed](1,-.5)--(1,.5) (5,-.5)--(5,.5);
 \draw[->,DeepBlue,thick](1,0)--(2.2,0)node[above]{$v_1$};\draw[->,DeepBlue,thick](5,0)--(6.2,0)node[above]{$v_2$};
 \node[below]at(3.5,-.8){控制体：等截面直管的两截面之间};
}{一维稳态模型中的相同截面积与流量；不是一般三维速度场的结论。}{ch13:quality:uniform-a}
图~\ref{ch13:quality:uniform-a}的侧壁不漏流，任取两端均有相同的截面积 $S$。若 $v_2>v_1$，每秒流出体积将大于流入体积，控制体内就不能保持处处充满且密度不变，与题设矛盾。
}
\newcommand{\ThAFallingFlow}{\ThermalFigure{
 \draw[DeepBlue,thick](-.8,0)..controls(-.5,-1)and(-.4,-2)..(-.35,-3);
 \draw[DeepBlue,thick](.8,0)..controls(.5,-1)and(.4,-2)..(.35,-3);
 \draw[thick](-.8,0)--(.8,0) (-.38,-2.5)--(.38,-2.5);
 \draw[<->](-1.6,0)--node[left]{$h$}(-1.6,-2.5);
 \draw[->](1.3,0)--(1.3,-.7)node[right]{$v$};\draw[->](1.3,-2.5)--(1.3,-3.6)node[right]{$v(h)$};
 \node[right]at(2.7,-1.2){$S\,v=S(h)v(h)$};
}{先用同一水微元的自由下落求速度，再用稳态连续性比较两个截面。}{ch13:quality:falling-a}
图~\ref{ch13:quality:falling-a}将两个步骤分开：高度差 $h$ 决定微元增加的动能，两个横截面则用于守恒单位时间的流体体积。流速增加时截面必须缩小，不能把两处面积也设为相同。
}
\newcommand{\ThAFlowMass}{\ThermalFigure{\ThFlowElement{
 \node[right]at(7,1){$\Delta m_A=\rho_AS_Av_A\Delta t$};
 \node[right]at(7,-.1){$\Delta m_B=\rho_BS_Bv_B\Delta t$};
}}{可压缩流的两体积柱可以不等大，但通过的质量须相同。}{ch13:quality:mass-a}
图~\ref{ch13:quality:mass-a}中每一柱体还须乘当地密度才得到质量。控制体内部的质量既不累积也不凭空产生，因此应令 $\Delta m_A=\Delta m_B$；只有同密度时才可进一步令两柱体积相等。
}
\newcommand{\ThDilutingBase}[1]{
 \draw[thick](-.6,0)--(-.6,-3) (.6,0)--(.6,-3);
 \draw[dashed](-1,0)--(1,0) (-1,-2.5)--(1,-2.5);
 \node[above]at(0,.2){$S,\rho,v$};\node[below]at(0,-3.2){$S,\rho(h),v(h)$};
 \draw[<->](-1.4,0)--node[left]{$h$}(-1.4,-2.5);
 \foreach \y in {-.2,-.4,-.6,-1.4,-1.9,-2.5}{\fill[DeepBlue](0,\y)circle(1.5pt);}#1}
\newcommand{\ThQDilutingColumn}{\ThermalFigure{\ThDilutingBase{}}{题设人为固定横截面积；分子沿竖直方向自由下落。}{ch13:quality:diluting-q}}
\newcommand{\ThADilutingColumn}{\ThermalFigure{\ThDilutingBase{
 \draw[->](1.5,-.1)--(1.5,-.7)node[right]{$v$};\draw[->](1.5,-2)--(1.5,-3)node[right]{$v(h)$};
}}{下方流速更大而面积固定，分子间的纵向间距增大。}{ch13:quality:diluting-a}
图~\ref{ch13:quality:diluting-a}中粒子点的疏密表示题设模型的密度变化。两个截面共用同一个 $S$，连续性要求 $\rho(h)v(h)=\rho v$；这里降低的是密度，不能照搬上一题缩小截面的真实水束模型。
}
\newcommand{\ThLayerBase}[1]{
 \fill[DeepBlue!10](0,1.5)rectangle(3,3);\fill[DeepBlue!25](0,0)rectangle(3,1.5);
 \draw[thick](0,3.3)--(0,0)--(3,0)--(3,3.3);\draw[DeepBlue](0,3)--(3,3) (0,1.5)--(3,1.5);
 \node at(1.5,2.3){$\rho_1,h_1$};\node at(1.5,.75){$\rho_2,h_2$};
 \draw[->](3.6,1.5)--(3.6,.2)node[right]{$x$};#1}
\newcommand{\ThQLayers}{\ThermalFigure{\ThLayerBase{\node[right]at(4.5,2.5){壁宽 $b$ 垂直纸面};}}{界面是下层深度坐标的起点，外侧大气压在净压力中抵消。}{ch13:quality:layers-q}}
\newcommand{\ThALayers}{\ThermalFigure{\ThLayerBase{
 \draw[DeepBlue,thick](3,1.5)--(4.2,1.5)--(5,0)--(3,0);
 \foreach \y in {.2,.65,1.1}{\draw[->,DeepBlue](3,\y)--({5-.8*\y/1.5},\y);}
 \node[right]at(5.3,.8){梯形表压分布};
}}{下层侧壁表压等于常量项与线性项之和，合力矩须逐条积分。}{ch13:quality:layers-a}
图~\ref{ch13:quality:layers-a}中下层顶部已经有上层液柱贡献的表压，故压力图是梯形而非从零开始的三角形。深度为 $x$ 的窄条面积为 $b\dd x$，其关于界面的力矩为 $x\Delta p(x)b\dd x$，由此得到合力的作用点。
}
\newcommand{\ThFloatBase}[1]{
 \fill[DeepBlue!10](-1,-1.5)rectangle(4,.4);\draw[DeepBlue](-1,.4)--(4,.4);
 \draw[thick](.7,-.8)rectangle(2.3,1.8);\node at(1.5,1.3){$m,S$};
 \draw[<->](2.9,-.8)--node[right]{$h_0$}(2.9,.4);#1}
\newcommand{\ThQFloat}{\ThermalFigure{\ThFloatBase{}}{直立柱体的平衡浸深；小振动中顶面仍露出、底面仍浸没。}{ch13:quality:float-q}}
\newcommand{\ThAFloat}{\ThermalFigure{\ThFloatBase{
 \draw[dashed](.7,-1.1)--(2.3,-1.1);\draw[->](.2,-.8)--(.2,-1.1)node[left]{$x$};
 \draw[->,thick](1.5,.1)--(1.5,-1.5)node[below]{$mg$};
 \draw[->,DeepBlue,thick](1.5,.1)--(1.5,1)node[right]{$F_b$};
}}{向下位移增加排液体积 $Sx$，额外浮力朝向上方。}{ch13:quality:float-a}
图~\ref{ch13:quality:float-a}将 $x$ 取向下为正；额外排液体积为 $Sx$，因此浮力增量的有向分量是 $-\rho gSx$。先用平衡条件消去 $mg-\rho gSh_0$，才留下关于平衡位置的振动方程。
}
\newcommand{\ThQOutlet}{\ThermalFigure{
 \draw[thick](0,3)--(0,0)--(3,0)--(3,.8) (3,1.1)--(3,3);\draw[DeepBlue](0,2.5)--(3,2.5);
 \node[left]at(0,2.5){$A$};\node[right]at(3.1,.95){$B$};\draw[->](3.1,.95)--(4.8,.95)node[right]{$v$};
 \draw[<->](5.5,.95)--node[right]{$h$}(5.5,2.5);
}{自由液面 $A$ 与出口 $B$ 均接大气，容器横截面远大于孔口。}{ch13:quality:outlet-q}}
\newcommand{\ThJetBase}[1]{
 \draw[DeepBlue,thick](0,.3)--(4,.3) (0,-.3)--(4,-.3);\draw[very thick](4,-1.6)--(4,1.6);
 \draw[->,thick](.6,0)--(2,0)node[midway,above]{$v$};\node[below]at(1,-.6){$A,\rho$};#1}
\newcommand{\ThQJet}{\ThermalFigure{\ThJetBase{\node[right]at(4.5,1){静止大平板};}}{自由射流垂直撞击平板后沿板面流走。}{ch13:quality:jet-q}}
\newcommand{\ThQMovingPlate}{\ThermalFigure{\ThJetBase{
 \draw[->,thick](4.5,1)--(6,1)node[right]{$u$};\node[right]at(4.5,-.8){$0\le u<v$};
}}{平板沿来流方向匀速后退，截面流速 $v$ 仍相对地面给定。}{ch13:quality:moving-q}}
\newcommand{\ThAMovingPlate}{\ThermalFigure{
 \draw[dashed](-.5,-1)rectangle(4.5,1);\draw[very thick](3.5,-1.4)--(3.5,1.4);
 \draw[->,DeepBlue,thick](0,0)--(2,0)node[midway,above]{$v-u$};
 \draw[->,DeepBlue,thick](3.5,.1)--(3.5,1.4);\draw[->,DeepBlue,thick](3.5,-.1)--(3.5,-1.4);
 \node[right,align=left]at(5,0){板系控制体静止\\出流法向速度为零};
}{在匀速平板的惯性系中计算撞击质量率和法向动量变化。}{ch13:quality:moving-a}
图~\ref{ch13:quality:moving-a}的虚线控制体随板一起平移，水以相对速度 $v-u$ 穿过入口。因此流入质量率和每单位质量损失的法向动量都含 $v-u$，合力中出现它的平方；向地面系中的平板传功时再乘板速 $u$。
}
\newcommand{\ThQSettling}{\ThermalFigure{
 \draw[thick](-2,2)--(-2,-2)--(2,-2)--(2,2);\draw[DeepBlue](-2,1.6)--(2,1.6);
 \draw[thick](0,0)circle(.55);\draw[->](0,0)--node[above]{$R$}(.55,0);
 \draw[->,DeepBlue,thick](1,0)--(1,-1.2)node[below]{$v$};
 \node[right,align=left]at(3,0){球 $\rho_s$\\液体 $\rho_f,\eta$};
}{小球远离容器壁的缓慢下沉模型；终端状态由合力为零确定。}{ch13:quality:settling-q}}
\newcommand{\ThDropletBase}[1]{
 \draw[DeepBlue,thick](0,0)circle(.65);\draw[DeepBlue,thick](2,0)circle(.65);
 \node[below]at(0,-.8){$R$};\node[below]at(2,-.8){$R$};\draw[->,thick](3.2,0)--(4.2,0);
 \draw[DeepBlue,thick](5.7,0)circle(.82);\node[below]at(5.7,-1){#1};}
\newcommand{\ThQDroplets}{\ThermalFigure{\ThDropletBase{$R_f$ 待求}}{两滴同种液体合并为一滴，体积守恒而界面面积改变。}{ch13:quality:droplets-q}}
\newcommand{\ThADroplets}{\ThermalFigure{\ThDropletBase{$R_f=2^{1/3}R$}
 \node[above]at(1,1.3){$A_i=8\pi R^2$};\node[above]at(5.7,1.3){$A_f=4\pi R_f^2$};
}{先由体积相加求半径，再比较新旧界面面积。}{ch13:quality:droplets-a}
图~\ref{ch13:quality:droplets-a}中的两个初始界面都计入 $A_i$，末态只剩一个界面。界面能的减少为 $\gamma(A_i-A_f)$；这些能量转移到流动和内能，不能解释为总能量消失。
}
\newcommand{\ThACoupledMercury}{\ThermalFigure{
 \draw[thick](0,-.4)rectangle(7,.4);\draw[line width=3pt](1,-.4)--(1,.4) (6,-.4)--(6,.4);
 \draw[->,DeepBlue,thick](1,1)--(2,1)node[right]{$q_1$};\draw[->,DeepBlue,thick](6,1)--(5,1)node[left]{$q_2$};
 \node at(3.5,0){绝热空气 $\bar p+\delta p$};
 \draw[<->](1,-1)--node[below]{$\ell-q_1-q_2$}(6,-1);
}{抽出两个水银内侧界面的位移，以空气柱的压缩量建立耦合项。}{ch13:quality:coupled-a}
图~\ref{ch13:quality:coupled-a}把两个位移都取向空气内部为正，所以气柱缩短量为 $q_1+q_2$。同一个 $\delta p$ 同时作用于两段水银，但两段惯性分别由各自总长度 $L_1,L_2$ 决定，不能只用竖直高度 $h$ 作为质量长度。
}
\newcommand{\ThACouette}{\ThermalFigure{
 \draw[thick](0,0)circle(1);\draw[thick](0,0)circle(2.3);\draw[DeepBlue,dashed](0,0)circle(1.6);
 \draw[->](0,0)--node[above]{$r$}(1.6,0);\fill(1.6,0)circle(1.5pt);
 \draw[->,DeepBlue,thick](1.6,0)--(1.6,1.2)node[right]{$v_\theta=r\omega$};
 \node[right,align=left]at(3,0){圆柱面面积 $2\pi r l$\\切应力 $\tau_{r\theta}$\\力矩再乘力臂 $r$};
}{在半径 $r$ 的同轴圆柱面上传递黏性力矩，不能只令切应力恒定。}{ch13:quality:couette-a}
图~\ref{ch13:quality:couette-a}的虚线圆柱面传递的总力矩为 $2\pi l r^2\tau_{r\theta}$。稳态无角动量积累，要求该乘积与 $r$ 无关；柱坐标中的 $-v_\theta/r$ 项保证刚体转动 $v_\theta\propto r$ 时没有虚假的剪切应力。
}
\newcommand{\ThAHangar}{\ThermalFigure{
 \draw[thick](-2,0)arc(180:0:2);\draw[thick](-2,0)--(2,0);\fill(0,0)circle(1pt)node[below]{$O$};
 \draw[->](0,0)--(1.414,1.414)node[midway,below right]{$R$};
 \draw[->,DeepBlue,thick](1.414,1.414)--(2.15,2.15)node[right]{$\Delta p\,\dd A$};
 \draw[->](.65,0)arc(0:45:.65)node[right]{$\theta$};\node[left]at(-2,0){$A$};
}{曲面微元的净压力沿外法线；竖直分量另乘 $\sin\theta$。}{ch13:quality:hangar-a}
图~\ref{ch13:quality:hangar-a}中微元面积为 $\dd A=lR\dd\theta$。压力差先由门口 $A$ 的停滞压确定库内压强，再与曲面外压比较；左右对称使水平分量相消，竖直积分还应包含法线的投影因子 $\sin\theta$。
}
\newcommand{\ThAFunnel}{\ThermalFigure{
 \draw[thick](-2,3)--(-.12,0) (.12,0)--(2,3);\draw[DeepBlue](-1.35,2)--(1.35,2);
 \draw[dashed](0,0)--(0,3.2);\draw[<->](2.5,0)--node[right]{$h$}(2.5,2);
 \draw[<->](-1.35,2.5)--node[above]{$2Rh/H$}(1.35,2.5);
 \draw[->,DeepBlue](0,0)--(0,-.8)node[right]{$Q$};
}{尖锥近似下自由液面的半径随水深变化，积分时不能保持水面面积不变。}{ch13:quality:funnel-a}
图~\ref{ch13:quality:funnel-a}中的相似三角形给出自由液面半径 $Rh/H$；液面下降 $\dd h<0$ 排出的体积是 $-A(h)\dd h$。将它等于小孔流出体积 $Q\dd t$，即可把流速公式转为排液时间积分。
}
\newcommand{\ThAShore}{\ThermalFigure{
 \draw[thick](-1,0)--(6,0)node[right]{岸线 $y=0$};\draw[->](0,0)--(0,3.4)node[above]{$y$};
 \draw[DeepBlue,thick](0,3)..controls(3,3)and(4,1.5)..(4.4,.3);\fill(0,3)circle(1.5pt)node[left]{$B$};
 \fill(3.2,1.9)circle(1.5pt)node[above right]{$P$};\draw[dashed](3.2,1.9)--(3.2,.6);
 \draw[->,thick](3.2,1.9)--(3.9,1.1);\draw(3.2,1.3)arc(270:311:.6)node[below]{$\theta$};
}{局部传播方向与朝岸法线 $-y$ 的夹角为 $\theta$；图示只标方向，不代替精确摆线。}{ch13:quality:shore-a}
图~\ref{ch13:quality:shore-a}中射线朝岸传播，所以 $\dd y<0$ 而 $\dd x>0$，这决定 $\dd x/\dd y=-\tan\theta$ 的负号。沿岸平移对称给出 $\sin\theta/c(y)$ 不变，初点 $B$ 的水平切向又给出其常量值。
}
\newcommand{\ThAHose}{\ThermalFigure{
 \coordinate(O)at(0,0);\coordinate(E)at(-.8,-2.3);\coordinate(P)at(1.5,-3.1);
 \draw[very thick](O)--(E)--(P);\fill(O)circle(2pt)node[above]{$O$};\node[left]at(E){$E$};\node[below]at(P){$P$};
 \draw[->,DeepBlue,thick](P)--(2.7,-3.5)node[right]{$u\vect e_2$};
 \fill(-.025,-1.925)circle(1.5pt)node[left]{$C$};\draw[->,thick](-.025,-1.925)--(-.025,-3.5)node[below]{$Mg$};
 \draw[dashed](0,0)--(1.5,-3.1);\node[right]at(2,-1){关于 $O$ 比较重力矩与出流角动量通量};
}{管和管内水组成控制体；出口动量通量关于接头 $O$ 有非零力矩。}{ch13:quality:hose-a}
图~\ref{ch13:quality:hose-a}的 $C$ 表示管与管内水的共同质心，两段均匀且等长时它位于两段中点的平均位置。入口位于 $O$，关于 $O$ 的入口角动量为零；出口 $P$ 则必须保留位置矢量与出流动量的叉乘。
}



{
\graphicspath{{../../../sample/xes-original/kinematics-supplement/}{../../../sample/xes-original/kinematics-supplement/assets/}}
\section{流线、流管与质量流率}
\newcommand{\SphereStreamlinesDiagram}{%
\begin{center}\begin{minipage}{\linewidth}\centering\begin{tikzpicture}[scale=.85,>=Latex]
 \shade[ball color=TextGray!35] (0,0) circle (1.2);
 \draw[Blue,thick] plot[smooth] coordinates {(-4.000000,0.354765) (-3.955556,0.354930) (-3.911111,0.355102) (-3.866667,0.355283) (-3.822222,0.355472) (-3.777778,0.355670) (-3.733333,0.355878) (-3.688889,0.356097) (-3.644444,0.356326) (-3.600000,0.356568) (-3.555556,0.356822) (-3.511111,0.357089) (-3.466667,0.357371) (-3.422222,0.357668) (-3.377778,0.357981) (-3.333333,0.358312) (-3.288889,0.358662) (-3.244444,0.359032) (-3.200000,0.359424) (-3.155556,0.359839) (-3.111111,0.360279) (-3.066667,0.360747) (-3.022222,0.361244) (-2.977778,0.361772) (-2.933333,0.362335) (-2.888889,0.362935) (-2.844444,0.363575) (-2.800000,0.364259) (-2.755556,0.364992) (-2.711111,0.365776) (-2.666667,0.366618) (-2.622222,0.367522) (-2.577778,0.368496) (-2.533333,0.369545) (-2.488889,0.370677) (-2.444444,0.371902) (-2.400000,0.373228) (-2.355556,0.374668) (-2.311111,0.376235) (-2.266667,0.377943) (-2.222222,0.379809) (-2.177778,0.381853) (-2.133333,0.384097) (-2.088889,0.386568) (-2.044444,0.389297) (-2.000000,0.392320) (-1.955556,0.395681) (-1.911111,0.399429) (-1.866667,0.403626) (-1.822222,0.408346) (-1.777778,0.413675) (-1.733333,0.419722) (-1.688889,0.426615) (-1.644444,0.434514) (-1.600000,0.443613) (-1.555556,0.454151) (-1.511111,0.466414) (-1.466667,0.480748) (-1.422222,0.497548) (-1.377778,0.517251) (-1.333333,0.540285) (-1.288889,0.566991) (-1.244444,0.597501) (-1.200000,0.631616) (-1.155556,0.668746) (-1.111111,0.707979) (-1.066667,0.748256) (-1.022222,0.788579) (-0.977778,0.828132) (-0.933333,0.866331) (-0.888889,0.902794) (-0.844444,0.937303) (-0.800000,0.969748) (-0.755556,1.000093) (-0.711111,1.028348) (-0.666667,1.054547) (-0.622222,1.078741) (-0.577778,1.100985) (-0.533333,1.121337) (-0.488889,1.139853) (-0.444444,1.156584) (-0.400000,1.171580) (-0.355556,1.184883) (-0.311111,1.196533) (-0.266667,1.206563) (-0.222222,1.215001) (-0.177778,1.221873) (-0.133333,1.227198) (-0.088889,1.230991) (-0.044444,1.233262) (0.000000,1.234019) (0.044444,1.233262) (0.088889,1.230991) (0.133333,1.227198) (0.177778,1.221873) (0.222222,1.215001) (0.266667,1.206563) (0.311111,1.196533) (0.355556,1.184883) (0.400000,1.171580) (0.444444,1.156584) (0.488889,1.139853) (0.533333,1.121337) (0.577778,1.100985) (0.622222,1.078741) (0.666667,1.054547) (0.711111,1.028348) (0.755556,1.000093) (0.800000,0.969748) (0.844444,0.937303) (0.888889,0.902794) (0.933333,0.866331) (0.977778,0.828132) (1.022222,0.788579) (1.066667,0.748256) (1.111111,0.707979) (1.155556,0.668746) (1.200000,0.631616) (1.244444,0.597501) (1.288889,0.566991) (1.333333,0.540285) (1.377778,0.517251) (1.422222,0.497548) (1.466667,0.480748) (1.511111,0.466414) (1.555556,0.454151) (1.600000,0.443613) (1.644444,0.434514) (1.688889,0.426615) (1.733333,0.419722) (1.777778,0.413675) (1.822222,0.408346) (1.866667,0.403626) (1.911111,0.399429) (1.955556,0.395681) (2.000000,0.392320) (2.044444,0.389297) (2.088889,0.386568) (2.133333,0.384097) (2.177778,0.381853) (2.222222,0.379809) (2.266667,0.377943) (2.311111,0.376235) (2.355556,0.374668) (2.400000,0.373228) (2.444444,0.371902) (2.488889,0.370677) (2.533333,0.369545) (2.577778,0.368496) (2.622222,0.367522) (2.666667,0.366618) (2.711111,0.365776) (2.755556,0.364992) (2.800000,0.364259) (2.844444,0.363575) (2.888889,0.362935) (2.933333,0.362335) (2.977778,0.361772) (3.022222,0.361244) (3.066667,0.360747) (3.111111,0.360279) (3.155556,0.359839) (3.200000,0.359424) (3.244444,0.359032) (3.288889,0.358662) (3.333333,0.358312) (3.377778,0.357981) (3.422222,0.357668) (3.466667,0.357371) (3.511111,0.357089) (3.555556,0.356822) (3.600000,0.356568) (3.644444,0.356326) (3.688889,0.356097) (3.733333,0.355878) (3.777778,0.355670) (3.822222,0.355472) (3.866667,0.355283) (3.911111,0.355102) (3.955556,0.354930) (4.000000,0.354765)};
\draw[Blue,thick,->] (2.444444,0.371902)--(2.711111,0.365776);
\draw[Blue,thick] plot[smooth] coordinates {(-4.000000,-0.354765) (-3.955556,-0.354930) (-3.911111,-0.355102) (-3.866667,-0.355283) (-3.822222,-0.355472) (-3.777778,-0.355670) (-3.733333,-0.355878) (-3.688889,-0.356097) (-3.644444,-0.356326) (-3.600000,-0.356568) (-3.555556,-0.356822) (-3.511111,-0.357089) (-3.466667,-0.357371) (-3.422222,-0.357668) (-3.377778,-0.357981) (-3.333333,-0.358312) (-3.288889,-0.358662) (-3.244444,-0.359032) (-3.200000,-0.359424) (-3.155556,-0.359839) (-3.111111,-0.360279) (-3.066667,-0.360747) (-3.022222,-0.361244) (-2.977778,-0.361772) (-2.933333,-0.362335) (-2.888889,-0.362935) (-2.844444,-0.363575) (-2.800000,-0.364259) (-2.755556,-0.364992) (-2.711111,-0.365776) (-2.666667,-0.366618) (-2.622222,-0.367522) (-2.577778,-0.368496) (-2.533333,-0.369545) (-2.488889,-0.370677) (-2.444444,-0.371902) (-2.400000,-0.373228) (-2.355556,-0.374668) (-2.311111,-0.376235) (-2.266667,-0.377943) (-2.222222,-0.379809) (-2.177778,-0.381853) (-2.133333,-0.384097) (-2.088889,-0.386568) (-2.044444,-0.389297) (-2.000000,-0.392320) (-1.955556,-0.395681) (-1.911111,-0.399429) (-1.866667,-0.403626) (-1.822222,-0.408346) (-1.777778,-0.413675) (-1.733333,-0.419722) (-1.688889,-0.426615) (-1.644444,-0.434514) (-1.600000,-0.443613) (-1.555556,-0.454151) (-1.511111,-0.466414) (-1.466667,-0.480748) (-1.422222,-0.497548) (-1.377778,-0.517251) (-1.333333,-0.540285) (-1.288889,-0.566991) (-1.244444,-0.597501) (-1.200000,-0.631616) (-1.155556,-0.668746) (-1.111111,-0.707979) (-1.066667,-0.748256) (-1.022222,-0.788579) (-0.977778,-0.828132) (-0.933333,-0.866331) (-0.888889,-0.902794) (-0.844444,-0.937303) (-0.800000,-0.969748) (-0.755556,-1.000093) (-0.711111,-1.028348) (-0.666667,-1.054547) (-0.622222,-1.078741) (-0.577778,-1.100985) (-0.533333,-1.121337) (-0.488889,-1.139853) (-0.444444,-1.156584) (-0.400000,-1.171580) (-0.355556,-1.184883) (-0.311111,-1.196533) (-0.266667,-1.206563) (-0.222222,-1.215001) (-0.177778,-1.221873) (-0.133333,-1.227198) (-0.088889,-1.230991) (-0.044444,-1.233262) (0.000000,-1.234019) (0.044444,-1.233262) (0.088889,-1.230991) (0.133333,-1.227198) (0.177778,-1.221873) (0.222222,-1.215001) (0.266667,-1.206563) (0.311111,-1.196533) (0.355556,-1.184883) (0.400000,-1.171580) (0.444444,-1.156584) (0.488889,-1.139853) (0.533333,-1.121337) (0.577778,-1.100985) (0.622222,-1.078741) (0.666667,-1.054547) (0.711111,-1.028348) (0.755556,-1.000093) (0.800000,-0.969748) (0.844444,-0.937303) (0.888889,-0.902794) (0.933333,-0.866331) (0.977778,-0.828132) (1.022222,-0.788579) (1.066667,-0.748256) (1.111111,-0.707979) (1.155556,-0.668746) (1.200000,-0.631616) (1.244444,-0.597501) (1.288889,-0.566991) (1.333333,-0.540285) (1.377778,-0.517251) (1.422222,-0.497548) (1.466667,-0.480748) (1.511111,-0.466414) (1.555556,-0.454151) (1.600000,-0.443613) (1.644444,-0.434514) (1.688889,-0.426615) (1.733333,-0.419722) (1.777778,-0.413675) (1.822222,-0.408346) (1.866667,-0.403626) (1.911111,-0.399429) (1.955556,-0.395681) (2.000000,-0.392320) (2.044444,-0.389297) (2.088889,-0.386568) (2.133333,-0.384097) (2.177778,-0.381853) (2.222222,-0.379809) (2.266667,-0.377943) (2.311111,-0.376235) (2.355556,-0.374668) (2.400000,-0.373228) (2.444444,-0.371902) (2.488889,-0.370677) (2.533333,-0.369545) (2.577778,-0.368496) (2.622222,-0.367522) (2.666667,-0.366618) (2.711111,-0.365776) (2.755556,-0.364992) (2.800000,-0.364259) (2.844444,-0.363575) (2.888889,-0.362935) (2.933333,-0.362335) (2.977778,-0.361772) (3.022222,-0.361244) (3.066667,-0.360747) (3.111111,-0.360279) (3.155556,-0.359839) (3.200000,-0.359424) (3.244444,-0.359032) (3.288889,-0.358662) (3.333333,-0.358312) (3.377778,-0.357981) (3.422222,-0.357668) (3.466667,-0.357371) (3.511111,-0.357089) (3.555556,-0.356822) (3.600000,-0.356568) (3.644444,-0.356326) (3.688889,-0.356097) (3.733333,-0.355878) (3.777778,-0.355670) (3.822222,-0.355472) (3.866667,-0.355283) (3.911111,-0.355102) (3.955556,-0.354930) (4.000000,-0.354765)};
\draw[Blue,thick,->] (2.444444,-0.371902)--(2.711111,-0.365776);
\draw[Blue,thick] plot[smooth] coordinates {(-4.000000,0.759789) (-3.955556,0.760117) (-3.911111,0.760460) (-3.866667,0.760818) (-3.822222,0.761193) (-3.777778,0.761585) (-3.733333,0.761996) (-3.688889,0.762426) (-3.644444,0.762877) (-3.600000,0.763350) (-3.555556,0.763847) (-3.511111,0.764369) (-3.466667,0.764917) (-3.422222,0.765493) (-3.377778,0.766100) (-3.333333,0.766738) (-3.288889,0.767411) (-3.244444,0.768120) (-3.200000,0.768868) (-3.155556,0.769657) (-3.111111,0.770491) (-3.066667,0.771373) (-3.022222,0.772307) (-2.977778,0.773295) (-2.933333,0.774342) (-2.888889,0.775453) (-2.844444,0.776632) (-2.800000,0.777885) (-2.755556,0.779218) (-2.711111,0.780637) (-2.666667,0.782148) (-2.622222,0.783760) (-2.577778,0.785481) (-2.533333,0.787320) (-2.488889,0.789287) (-2.444444,0.791394) (-2.400000,0.793653) (-2.355556,0.796078) (-2.311111,0.798683) (-2.266667,0.801485) (-2.222222,0.804502) (-2.177778,0.807756) (-2.133333,0.811267) (-2.088889,0.815062) (-2.044444,0.819167) (-2.000000,0.823613) (-1.955556,0.828432) (-1.911111,0.833662) (-1.866667,0.839340) (-1.822222,0.845511) (-1.777778,0.852221) (-1.733333,0.859519) (-1.688889,0.867457) (-1.644444,0.876089) (-1.600000,0.885468) (-1.555556,0.895648) (-1.511111,0.906679) (-1.466667,0.918602) (-1.422222,0.931452) (-1.377778,0.945247) (-1.333333,0.959987) (-1.288889,0.975652) (-1.244444,0.992195) (-1.200000,1.009543) (-1.155556,1.027597) (-1.111111,1.046235) (-1.066667,1.065313) (-1.022222,1.084677) (-0.977778,1.104165) (-0.933333,1.123613) (-0.888889,1.142866) (-0.844444,1.161777) (-0.800000,1.180214) (-0.755556,1.198058) (-0.711111,1.215207) (-0.666667,1.231576) (-0.622222,1.247091) (-0.577778,1.261692) (-0.533333,1.275330) (-0.488889,1.287967) (-0.444444,1.299572) (-0.400000,1.310119) (-0.355556,1.319591) (-0.311111,1.327973) (-0.266667,1.335253) (-0.222222,1.341424) (-0.177778,1.346479) (-0.133333,1.350415) (-0.088889,1.353228) (-0.044444,1.354917) (0.000000,1.355480) (0.044444,1.354917) (0.088889,1.353228) (0.133333,1.350415) (0.177778,1.346479) (0.222222,1.341424) (0.266667,1.335253) (0.311111,1.327973) (0.355556,1.319591) (0.400000,1.310119) (0.444444,1.299572) (0.488889,1.287967) (0.533333,1.275330) (0.577778,1.261692) (0.622222,1.247091) (0.666667,1.231576) (0.711111,1.215207) (0.755556,1.198058) (0.800000,1.180214) (0.844444,1.161777) (0.888889,1.142866) (0.933333,1.123613) (0.977778,1.104165) (1.022222,1.084677) (1.066667,1.065313) (1.111111,1.046235) (1.155556,1.027597) (1.200000,1.009543) (1.244444,0.992195) (1.288889,0.975652) (1.333333,0.959987) (1.377778,0.945247) (1.422222,0.931452) (1.466667,0.918602) (1.511111,0.906679) (1.555556,0.895648) (1.600000,0.885468) (1.644444,0.876089) (1.688889,0.867457) (1.733333,0.859519) (1.777778,0.852221) (1.822222,0.845511) (1.866667,0.839340) (1.911111,0.833662) (1.955556,0.828432) (2.000000,0.823613) (2.044444,0.819167) (2.088889,0.815062) (2.133333,0.811267) (2.177778,0.807756) (2.222222,0.804502) (2.266667,0.801485) (2.311111,0.798683) (2.355556,0.796078) (2.400000,0.793653) (2.444444,0.791394) (2.488889,0.789287) (2.533333,0.787320) (2.577778,0.785481) (2.622222,0.783760) (2.666667,0.782148) (2.711111,0.780637) (2.755556,0.779218) (2.800000,0.777885) (2.844444,0.776632) (2.888889,0.775453) (2.933333,0.774342) (2.977778,0.773295) (3.022222,0.772307) (3.066667,0.771373) (3.111111,0.770491) (3.155556,0.769657) (3.200000,0.768868) (3.244444,0.768120) (3.288889,0.767411) (3.333333,0.766738) (3.377778,0.766100) (3.422222,0.765493) (3.466667,0.764917) (3.511111,0.764369) (3.555556,0.763847) (3.600000,0.763350) (3.644444,0.762877) (3.688889,0.762426) (3.733333,0.761996) (3.777778,0.761585) (3.822222,0.761193) (3.866667,0.760818) (3.911111,0.760460) (3.955556,0.760117) (4.000000,0.759789)};
\draw[Blue,thick,->] (2.444444,0.791394)--(2.711111,0.780637);
\draw[Blue,thick] plot[smooth] coordinates {(-4.000000,-0.759789) (-3.955556,-0.760117) (-3.911111,-0.760460) (-3.866667,-0.760818) (-3.822222,-0.761193) (-3.777778,-0.761585) (-3.733333,-0.761996) (-3.688889,-0.762426) (-3.644444,-0.762877) (-3.600000,-0.763350) (-3.555556,-0.763847) (-3.511111,-0.764369) (-3.466667,-0.764917) (-3.422222,-0.765493) (-3.377778,-0.766100) (-3.333333,-0.766738) (-3.288889,-0.767411) (-3.244444,-0.768120) (-3.200000,-0.768868) (-3.155556,-0.769657) (-3.111111,-0.770491) (-3.066667,-0.771373) (-3.022222,-0.772307) (-2.977778,-0.773295) (-2.933333,-0.774342) (-2.888889,-0.775453) (-2.844444,-0.776632) (-2.800000,-0.777885) (-2.755556,-0.779218) (-2.711111,-0.780637) (-2.666667,-0.782148) (-2.622222,-0.783760) (-2.577778,-0.785481) (-2.533333,-0.787320) (-2.488889,-0.789287) (-2.444444,-0.791394) (-2.400000,-0.793653) (-2.355556,-0.796078) (-2.311111,-0.798683) (-2.266667,-0.801485) (-2.222222,-0.804502) (-2.177778,-0.807756) (-2.133333,-0.811267) (-2.088889,-0.815062) (-2.044444,-0.819167) (-2.000000,-0.823613) (-1.955556,-0.828432) (-1.911111,-0.833662) (-1.866667,-0.839340) (-1.822222,-0.845511) (-1.777778,-0.852221) (-1.733333,-0.859519) (-1.688889,-0.867457) (-1.644444,-0.876089) (-1.600000,-0.885468) (-1.555556,-0.895648) (-1.511111,-0.906679) (-1.466667,-0.918602) (-1.422222,-0.931452) (-1.377778,-0.945247) (-1.333333,-0.959987) (-1.288889,-0.975652) (-1.244444,-0.992195) (-1.200000,-1.009543) (-1.155556,-1.027597) (-1.111111,-1.046235) (-1.066667,-1.065313) (-1.022222,-1.084677) (-0.977778,-1.104165) (-0.933333,-1.123613) (-0.888889,-1.142866) (-0.844444,-1.161777) (-0.800000,-1.180214) (-0.755556,-1.198058) (-0.711111,-1.215207) (-0.666667,-1.231576) (-0.622222,-1.247091) (-0.577778,-1.261692) (-0.533333,-1.275330) (-0.488889,-1.287967) (-0.444444,-1.299572) (-0.400000,-1.310119) (-0.355556,-1.319591) (-0.311111,-1.327973) (-0.266667,-1.335253) (-0.222222,-1.341424) (-0.177778,-1.346479) (-0.133333,-1.350415) (-0.088889,-1.353228) (-0.044444,-1.354917) (0.000000,-1.355480) (0.044444,-1.354917) (0.088889,-1.353228) (0.133333,-1.350415) (0.177778,-1.346479) (0.222222,-1.341424) (0.266667,-1.335253) (0.311111,-1.327973) (0.355556,-1.319591) (0.400000,-1.310119) (0.444444,-1.299572) (0.488889,-1.287967) (0.533333,-1.275330) (0.577778,-1.261692) (0.622222,-1.247091) (0.666667,-1.231576) (0.711111,-1.215207) (0.755556,-1.198058) (0.800000,-1.180214) (0.844444,-1.161777) (0.888889,-1.142866) (0.933333,-1.123613) (0.977778,-1.104165) (1.022222,-1.084677) (1.066667,-1.065313) (1.111111,-1.046235) (1.155556,-1.027597) (1.200000,-1.009543) (1.244444,-0.992195) (1.288889,-0.975652) (1.333333,-0.959987) (1.377778,-0.945247) (1.422222,-0.931452) (1.466667,-0.918602) (1.511111,-0.906679) (1.555556,-0.895648) (1.600000,-0.885468) (1.644444,-0.876089) (1.688889,-0.867457) (1.733333,-0.859519) (1.777778,-0.852221) (1.822222,-0.845511) (1.866667,-0.839340) (1.911111,-0.833662) (1.955556,-0.828432) (2.000000,-0.823613) (2.044444,-0.819167) (2.088889,-0.815062) (2.133333,-0.811267) (2.177778,-0.807756) (2.222222,-0.804502) (2.266667,-0.801485) (2.311111,-0.798683) (2.355556,-0.796078) (2.400000,-0.793653) (2.444444,-0.791394) (2.488889,-0.789287) (2.533333,-0.787320) (2.577778,-0.785481) (2.622222,-0.783760) (2.666667,-0.782148) (2.711111,-0.780637) (2.755556,-0.779218) (2.800000,-0.777885) (2.844444,-0.776632) (2.888889,-0.775453) (2.933333,-0.774342) (2.977778,-0.773295) (3.022222,-0.772307) (3.066667,-0.771373) (3.111111,-0.770491) (3.155556,-0.769657) (3.200000,-0.768868) (3.244444,-0.768120) (3.288889,-0.767411) (3.333333,-0.766738) (3.377778,-0.766100) (3.422222,-0.765493) (3.466667,-0.764917) (3.511111,-0.764369) (3.555556,-0.763847) (3.600000,-0.763350) (3.644444,-0.762877) (3.688889,-0.762426) (3.733333,-0.761996) (3.777778,-0.761585) (3.822222,-0.761193) (3.866667,-0.760818) (3.911111,-0.760460) (3.955556,-0.760117) (4.000000,-0.759789)};
\draw[Blue,thick,->] (2.444444,-0.791394)--(2.711111,-0.780637);
\draw[Blue,thick] plot[smooth] coordinates {(-4.000000,1.214450) (-3.955556,1.214907) (-3.911111,1.215383) (-3.866667,1.215879) (-3.822222,1.216396) (-3.777778,1.216935) (-3.733333,1.217498) (-3.688889,1.218085) (-3.644444,1.218699) (-3.600000,1.219339) (-3.555556,1.220009) (-3.511111,1.220709) (-3.466667,1.221441) (-3.422222,1.222207) (-3.377778,1.223009) (-3.333333,1.223849) (-3.288889,1.224730) (-3.244444,1.225652) (-3.200000,1.226620) (-3.155556,1.227635) (-3.111111,1.228701) (-3.066667,1.229820) (-3.022222,1.230995) (-2.977778,1.232231) (-2.933333,1.233531) (-2.888889,1.234898) (-2.844444,1.236338) (-2.800000,1.237854) (-2.755556,1.239450) (-2.711111,1.241134) (-2.666667,1.242909) (-2.622222,1.244781) (-2.577778,1.246757) (-2.533333,1.248844) (-2.488889,1.251047) (-2.444444,1.253376) (-2.400000,1.255837) (-2.355556,1.258439) (-2.311111,1.261191) (-2.266667,1.264102) (-2.222222,1.267182) (-2.177778,1.270441) (-2.133333,1.273890) (-2.088889,1.277540) (-2.044444,1.281403) (-2.000000,1.285489) (-1.955556,1.289811) (-1.911111,1.294382) (-1.866667,1.299212) (-1.822222,1.304313) (-1.777778,1.309697) (-1.733333,1.315374) (-1.688889,1.321353) (-1.644444,1.327641) (-1.600000,1.334246) (-1.555556,1.341170) (-1.511111,1.348415) (-1.466667,1.355979) (-1.422222,1.363858) (-1.377778,1.372040) (-1.333333,1.380514) (-1.288889,1.389261) (-1.244444,1.398259) (-1.200000,1.407481) (-1.155556,1.416894) (-1.111111,1.426464) (-1.066667,1.436149) (-1.022222,1.445908) (-0.977778,1.455693) (-0.933333,1.465457) (-0.888889,1.475150) (-0.844444,1.484722) (-0.800000,1.494124) (-0.755556,1.503306) (-0.711111,1.512221) (-0.666667,1.520822) (-0.622222,1.529066) (-0.577778,1.536914) (-0.533333,1.544325) (-0.488889,1.551267) (-0.444444,1.557706) (-0.400000,1.563615) (-0.355556,1.568967) (-0.311111,1.573741) (-0.266667,1.577916) (-0.222222,1.581477) (-0.177778,1.584409) (-0.133333,1.586700) (-0.088889,1.588343) (-0.044444,1.589332) (0.000000,1.589662) (0.044444,1.589332) (0.088889,1.588343) (0.133333,1.586700) (0.177778,1.584409) (0.222222,1.581477) (0.266667,1.577916) (0.311111,1.573741) (0.355556,1.568967) (0.400000,1.563615) (0.444444,1.557706) (0.488889,1.551267) (0.533333,1.544325) (0.577778,1.536914) (0.622222,1.529066) (0.666667,1.520822) (0.711111,1.512221) (0.755556,1.503306) (0.800000,1.494124) (0.844444,1.484722) (0.888889,1.475150) (0.933333,1.465457) (0.977778,1.455693) (1.022222,1.445908) (1.066667,1.436149) (1.111111,1.426464) (1.155556,1.416894) (1.200000,1.407481) (1.244444,1.398259) (1.288889,1.389261) (1.333333,1.380514) (1.377778,1.372040) (1.422222,1.363858) (1.466667,1.355979) (1.511111,1.348415) (1.555556,1.341170) (1.600000,1.334246) (1.644444,1.327641) (1.688889,1.321353) (1.733333,1.315374) (1.777778,1.309697) (1.822222,1.304313) (1.866667,1.299212) (1.911111,1.294382) (1.955556,1.289811) (2.000000,1.285489) (2.044444,1.281403) (2.088889,1.277540) (2.133333,1.273890) (2.177778,1.270441) (2.222222,1.267182) (2.266667,1.264102) (2.311111,1.261191) (2.355556,1.258439) (2.400000,1.255837) (2.444444,1.253376) (2.488889,1.251047) (2.533333,1.248844) (2.577778,1.246757) (2.622222,1.244781) (2.666667,1.242909) (2.711111,1.241134) (2.755556,1.239450) (2.800000,1.237854) (2.844444,1.236338) (2.888889,1.234898) (2.933333,1.233531) (2.977778,1.232231) (3.022222,1.230995) (3.066667,1.229820) (3.111111,1.228701) (3.155556,1.227635) (3.200000,1.226620) (3.244444,1.225652) (3.288889,1.224730) (3.333333,1.223849) (3.377778,1.223009) (3.422222,1.222207) (3.466667,1.221441) (3.511111,1.220709) (3.555556,1.220009) (3.600000,1.219339) (3.644444,1.218699) (3.688889,1.218085) (3.733333,1.217498) (3.777778,1.216935) (3.822222,1.216396) (3.866667,1.215879) (3.911111,1.215383) (3.955556,1.214907) (4.000000,1.214450)};
\draw[Blue,thick,->] (2.444444,1.253376)--(2.711111,1.241134);
\draw[Blue,thick] plot[smooth] coordinates {(-4.000000,-1.214450) (-3.955556,-1.214907) (-3.911111,-1.215383) (-3.866667,-1.215879) (-3.822222,-1.216396) (-3.777778,-1.216935) (-3.733333,-1.217498) (-3.688889,-1.218085) (-3.644444,-1.218699) (-3.600000,-1.219339) (-3.555556,-1.220009) (-3.511111,-1.220709) (-3.466667,-1.221441) (-3.422222,-1.222207) (-3.377778,-1.223009) (-3.333333,-1.223849) (-3.288889,-1.224730) (-3.244444,-1.225652) (-3.200000,-1.226620) (-3.155556,-1.227635) (-3.111111,-1.228701) (-3.066667,-1.229820) (-3.022222,-1.230995) (-2.977778,-1.232231) (-2.933333,-1.233531) (-2.888889,-1.234898) (-2.844444,-1.236338) (-2.800000,-1.237854) (-2.755556,-1.239450) (-2.711111,-1.241134) (-2.666667,-1.242909) (-2.622222,-1.244781) (-2.577778,-1.246757) (-2.533333,-1.248844) (-2.488889,-1.251047) (-2.444444,-1.253376) (-2.400000,-1.255837) (-2.355556,-1.258439) (-2.311111,-1.261191) (-2.266667,-1.264102) (-2.222222,-1.267182) (-2.177778,-1.270441) (-2.133333,-1.273890) (-2.088889,-1.277540) (-2.044444,-1.281403) (-2.000000,-1.285489) (-1.955556,-1.289811) (-1.911111,-1.294382) (-1.866667,-1.299212) (-1.822222,-1.304313) (-1.777778,-1.309697) (-1.733333,-1.315374) (-1.688889,-1.321353) (-1.644444,-1.327641) (-1.600000,-1.334246) (-1.555556,-1.341170) (-1.511111,-1.348415) (-1.466667,-1.355979) (-1.422222,-1.363858) (-1.377778,-1.372040) (-1.333333,-1.380514) (-1.288889,-1.389261) (-1.244444,-1.398259) (-1.200000,-1.407481) (-1.155556,-1.416894) (-1.111111,-1.426464) (-1.066667,-1.436149) (-1.022222,-1.445908) (-0.977778,-1.455693) (-0.933333,-1.465457) (-0.888889,-1.475150) (-0.844444,-1.484722) (-0.800000,-1.494124) (-0.755556,-1.503306) (-0.711111,-1.512221) (-0.666667,-1.520822) (-0.622222,-1.529066) (-0.577778,-1.536914) (-0.533333,-1.544325) (-0.488889,-1.551267) (-0.444444,-1.557706) (-0.400000,-1.563615) (-0.355556,-1.568967) (-0.311111,-1.573741) (-0.266667,-1.577916) (-0.222222,-1.581477) (-0.177778,-1.584409) (-0.133333,-1.586700) (-0.088889,-1.588343) (-0.044444,-1.589332) (0.000000,-1.589662) (0.044444,-1.589332) (0.088889,-1.588343) (0.133333,-1.586700) (0.177778,-1.584409) (0.222222,-1.581477) (0.266667,-1.577916) (0.311111,-1.573741) (0.355556,-1.568967) (0.400000,-1.563615) (0.444444,-1.557706) (0.488889,-1.551267) (0.533333,-1.544325) (0.577778,-1.536914) (0.622222,-1.529066) (0.666667,-1.520822) (0.711111,-1.512221) (0.755556,-1.503306) (0.800000,-1.494124) (0.844444,-1.484722) (0.888889,-1.475150) (0.933333,-1.465457) (0.977778,-1.455693) (1.022222,-1.445908) (1.066667,-1.436149) (1.111111,-1.426464) (1.155556,-1.416894) (1.200000,-1.407481) (1.244444,-1.398259) (1.288889,-1.389261) (1.333333,-1.380514) (1.377778,-1.372040) (1.422222,-1.363858) (1.466667,-1.355979) (1.511111,-1.348415) (1.555556,-1.341170) (1.600000,-1.334246) (1.644444,-1.327641) (1.688889,-1.321353) (1.733333,-1.315374) (1.777778,-1.309697) (1.822222,-1.304313) (1.866667,-1.299212) (1.911111,-1.294382) (1.955556,-1.289811) (2.000000,-1.285489) (2.044444,-1.281403) (2.088889,-1.277540) (2.133333,-1.273890) (2.177778,-1.270441) (2.222222,-1.267182) (2.266667,-1.264102) (2.311111,-1.261191) (2.355556,-1.258439) (2.400000,-1.255837) (2.444444,-1.253376) (2.488889,-1.251047) (2.533333,-1.248844) (2.577778,-1.246757) (2.622222,-1.244781) (2.666667,-1.242909) (2.711111,-1.241134) (2.755556,-1.239450) (2.800000,-1.237854) (2.844444,-1.236338) (2.888889,-1.234898) (2.933333,-1.233531) (2.977778,-1.232231) (3.022222,-1.230995) (3.066667,-1.229820) (3.111111,-1.228701) (3.155556,-1.227635) (3.200000,-1.226620) (3.244444,-1.225652) (3.288889,-1.224730) (3.333333,-1.223849) (3.377778,-1.223009) (3.422222,-1.222207) (3.466667,-1.221441) (3.511111,-1.220709) (3.555556,-1.220009) (3.600000,-1.219339) (3.644444,-1.218699) (3.688889,-1.218085) (3.733333,-1.217498) (3.777778,-1.216935) (3.822222,-1.216396) (3.866667,-1.215879) (3.911111,-1.215383) (3.955556,-1.214907) (4.000000,-1.214450)};
\draw[Blue,thick,->] (2.444444,-1.253376)--(2.711111,-1.241134);
\draw[Blue,thick] plot[smooth] coordinates {(-4.000000,1.768360) (-3.955556,1.768887) (-3.911111,1.769433) (-3.866667,1.769999) (-3.822222,1.770586) (-3.777778,1.771195) (-3.733333,1.771827) (-3.688889,1.772483) (-3.644444,1.773164) (-3.600000,1.773870) (-3.555556,1.774604) (-3.511111,1.775367) (-3.466667,1.776159) (-3.422222,1.776982) (-3.377778,1.777837) (-3.333333,1.778726) (-3.288889,1.779651) (-3.244444,1.780612) (-3.200000,1.781612) (-3.155556,1.782652) (-3.111111,1.783734) (-3.066667,1.784860) (-3.022222,1.786032) (-2.977778,1.787252) (-2.933333,1.788522) (-2.888889,1.789844) (-2.844444,1.791220) (-2.800000,1.792653) (-2.755556,1.794145) (-2.711111,1.795699) (-2.666667,1.797317) (-2.622222,1.799002) (-2.577778,1.800757) (-2.533333,1.802583) (-2.488889,1.804485) (-2.444444,1.806464) (-2.400000,1.808524) (-2.355556,1.810668) (-2.311111,1.812898) (-2.266667,1.815217) (-2.222222,1.817627) (-2.177778,1.820132) (-2.133333,1.822734) (-2.088889,1.825436) (-2.044444,1.828239) (-2.000000,1.831145) (-1.955556,1.834157) (-1.911111,1.837274) (-1.866667,1.840500) (-1.822222,1.843833) (-1.777778,1.847274) (-1.733333,1.850822) (-1.688889,1.854477) (-1.644444,1.858236) (-1.600000,1.862097) (-1.555556,1.866056) (-1.511111,1.870110) (-1.466667,1.874252) (-1.422222,1.878478) (-1.377778,1.882779) (-1.333333,1.887149) (-1.288889,1.891577) (-1.244444,1.896054) (-1.200000,1.900569) (-1.155556,1.905109) (-1.111111,1.909662) (-1.066667,1.914215) (-1.022222,1.918751) (-0.977778,1.923258) (-0.933333,1.927717) (-0.888889,1.932115) (-0.844444,1.936433) (-0.800000,1.940655) (-0.755556,1.944765) (-0.711111,1.948746) (-0.666667,1.952581) (-0.622222,1.956254) (-0.577778,1.959750) (-0.533333,1.963053) (-0.488889,1.966149) (-0.444444,1.969024) (-0.400000,1.971666) (-0.355556,1.974062) (-0.311111,1.976203) (-0.266667,1.978078) (-0.222222,1.979679) (-0.177778,1.980998) (-0.133333,1.982031) (-0.088889,1.982772) (-0.044444,1.983218) (0.000000,1.983367) (0.044444,1.983218) (0.088889,1.982772) (0.133333,1.982031) (0.177778,1.980998) (0.222222,1.979679) (0.266667,1.978078) (0.311111,1.976203) (0.355556,1.974062) (0.400000,1.971666) (0.444444,1.969024) (0.488889,1.966149) (0.533333,1.963053) (0.577778,1.959750) (0.622222,1.956254) (0.666667,1.952581) (0.711111,1.948746) (0.755556,1.944765) (0.800000,1.940655) (0.844444,1.936433) (0.888889,1.932115) (0.933333,1.927717) (0.977778,1.923258) (1.022222,1.918751) (1.066667,1.914215) (1.111111,1.909662) (1.155556,1.905109) (1.200000,1.900569) (1.244444,1.896054) (1.288889,1.891577) (1.333333,1.887149) (1.377778,1.882779) (1.422222,1.878478) (1.466667,1.874252) (1.511111,1.870110) (1.555556,1.866056) (1.600000,1.862097) (1.644444,1.858236) (1.688889,1.854477) (1.733333,1.850822) (1.777778,1.847274) (1.822222,1.843833) (1.866667,1.840500) (1.911111,1.837274) (1.955556,1.834157) (2.000000,1.831145) (2.044444,1.828239) (2.088889,1.825436) (2.133333,1.822734) (2.177778,1.820132) (2.222222,1.817627) (2.266667,1.815217) (2.311111,1.812898) (2.355556,1.810668) (2.400000,1.808524) (2.444444,1.806464) (2.488889,1.804485) (2.533333,1.802583) (2.577778,1.800757) (2.622222,1.799002) (2.666667,1.797317) (2.711111,1.795699) (2.755556,1.794145) (2.800000,1.792653) (2.844444,1.791220) (2.888889,1.789844) (2.933333,1.788522) (2.977778,1.787252) (3.022222,1.786032) (3.066667,1.784860) (3.111111,1.783734) (3.155556,1.782652) (3.200000,1.781612) (3.244444,1.780612) (3.288889,1.779651) (3.333333,1.778726) (3.377778,1.777837) (3.422222,1.776982) (3.466667,1.776159) (3.511111,1.775367) (3.555556,1.774604) (3.600000,1.773870) (3.644444,1.773164) (3.688889,1.772483) (3.733333,1.771827) (3.777778,1.771195) (3.822222,1.770586) (3.866667,1.769999) (3.911111,1.769433) (3.955556,1.768887) (4.000000,1.768360)};
\draw[Blue,thick,->] (2.444444,1.806464)--(2.711111,1.795699);
\draw[Blue,thick] plot[smooth] coordinates {(-4.000000,-1.768360) (-3.955556,-1.768887) (-3.911111,-1.769433) (-3.866667,-1.769999) (-3.822222,-1.770586) (-3.777778,-1.771195) (-3.733333,-1.771827) (-3.688889,-1.772483) (-3.644444,-1.773164) (-3.600000,-1.773870) (-3.555556,-1.774604) (-3.511111,-1.775367) (-3.466667,-1.776159) (-3.422222,-1.776982) (-3.377778,-1.777837) (-3.333333,-1.778726) (-3.288889,-1.779651) (-3.244444,-1.780612) (-3.200000,-1.781612) (-3.155556,-1.782652) (-3.111111,-1.783734) (-3.066667,-1.784860) (-3.022222,-1.786032) (-2.977778,-1.787252) (-2.933333,-1.788522) (-2.888889,-1.789844) (-2.844444,-1.791220) (-2.800000,-1.792653) (-2.755556,-1.794145) (-2.711111,-1.795699) (-2.666667,-1.797317) (-2.622222,-1.799002) (-2.577778,-1.800757) (-2.533333,-1.802583) (-2.488889,-1.804485) (-2.444444,-1.806464) (-2.400000,-1.808524) (-2.355556,-1.810668) (-2.311111,-1.812898) (-2.266667,-1.815217) (-2.222222,-1.817627) (-2.177778,-1.820132) (-2.133333,-1.822734) (-2.088889,-1.825436) (-2.044444,-1.828239) (-2.000000,-1.831145) (-1.955556,-1.834157) (-1.911111,-1.837274) (-1.866667,-1.840500) (-1.822222,-1.843833) (-1.777778,-1.847274) (-1.733333,-1.850822) (-1.688889,-1.854477) (-1.644444,-1.858236) (-1.600000,-1.862097) (-1.555556,-1.866056) (-1.511111,-1.870110) (-1.466667,-1.874252) (-1.422222,-1.878478) (-1.377778,-1.882779) (-1.333333,-1.887149) (-1.288889,-1.891577) (-1.244444,-1.896054) (-1.200000,-1.900569) (-1.155556,-1.905109) (-1.111111,-1.909662) (-1.066667,-1.914215) (-1.022222,-1.918751) (-0.977778,-1.923258) (-0.933333,-1.927717) (-0.888889,-1.932115) (-0.844444,-1.936433) (-0.800000,-1.940655) (-0.755556,-1.944765) (-0.711111,-1.948746) (-0.666667,-1.952581) (-0.622222,-1.956254) (-0.577778,-1.959750) (-0.533333,-1.963053) (-0.488889,-1.966149) (-0.444444,-1.969024) (-0.400000,-1.971666) (-0.355556,-1.974062) (-0.311111,-1.976203) (-0.266667,-1.978078) (-0.222222,-1.979679) (-0.177778,-1.980998) (-0.133333,-1.982031) (-0.088889,-1.982772) (-0.044444,-1.983218) (0.000000,-1.983367) (0.044444,-1.983218) (0.088889,-1.982772) (0.133333,-1.982031) (0.177778,-1.980998) (0.222222,-1.979679) (0.266667,-1.978078) (0.311111,-1.976203) (0.355556,-1.974062) (0.400000,-1.971666) (0.444444,-1.969024) (0.488889,-1.966149) (0.533333,-1.963053) (0.577778,-1.959750) (0.622222,-1.956254) (0.666667,-1.952581) (0.711111,-1.948746) (0.755556,-1.944765) (0.800000,-1.940655) (0.844444,-1.936433) (0.888889,-1.932115) (0.933333,-1.927717) (0.977778,-1.923258) (1.022222,-1.918751) (1.066667,-1.914215) (1.111111,-1.909662) (1.155556,-1.905109) (1.200000,-1.900569) (1.244444,-1.896054) (1.288889,-1.891577) (1.333333,-1.887149) (1.377778,-1.882779) (1.422222,-1.878478) (1.466667,-1.874252) (1.511111,-1.870110) (1.555556,-1.866056) (1.600000,-1.862097) (1.644444,-1.858236) (1.688889,-1.854477) (1.733333,-1.850822) (1.777778,-1.847274) (1.822222,-1.843833) (1.866667,-1.840500) (1.911111,-1.837274) (1.955556,-1.834157) (2.000000,-1.831145) (2.044444,-1.828239) (2.088889,-1.825436) (2.133333,-1.822734) (2.177778,-1.820132) (2.222222,-1.817627) (2.266667,-1.815217) (2.311111,-1.812898) (2.355556,-1.810668) (2.400000,-1.808524) (2.444444,-1.806464) (2.488889,-1.804485) (2.533333,-1.802583) (2.577778,-1.800757) (2.622222,-1.799002) (2.666667,-1.797317) (2.711111,-1.795699) (2.755556,-1.794145) (2.800000,-1.792653) (2.844444,-1.791220) (2.888889,-1.789844) (2.933333,-1.788522) (2.977778,-1.787252) (3.022222,-1.786032) (3.066667,-1.784860) (3.111111,-1.783734) (3.155556,-1.782652) (3.200000,-1.781612) (3.244444,-1.780612) (3.288889,-1.779651) (3.333333,-1.778726) (3.377778,-1.777837) (3.422222,-1.776982) (3.466667,-1.776159) (3.511111,-1.775367) (3.555556,-1.774604) (3.600000,-1.773870) (3.644444,-1.773164) (3.688889,-1.772483) (3.733333,-1.771827) (3.777778,-1.771195) (3.822222,-1.770586) (3.866667,-1.769999) (3.911111,-1.769433) (3.955556,-1.768887) (4.000000,-1.768360)};
\draw[Blue,thick,->] (2.444444,-1.806464)--(2.711111,-1.795699);

 \draw[Blue,thick,->] (-4,0)--(-1.25,0);\draw[Blue,thick,->] (1.25,0)--(4,0);
\end{tikzpicture}\figcap{圆球绕流的子午面流线。不同流线绕过球面，不在球顶或球底合并。}\label{ch13:fig:k10-2-23}\end{minipage}\end{center}}

\newcommand{\StreamtubeDiagram}{%
\begin{center}\begin{minipage}{\linewidth}\centering\begin{tikzpicture}[scale=.95,>=Latex]
 % Center c(t)=(t,0,.1t^2); section basis e_y and (-.2t,0,1)/sqrt(1+.04t^2).
 \foreach \ang in {0,60,120,180,240,300}{
  \draw[Blue,thick,domain=0:5,samples=75,variable=\t]
   plot ({\t+(.70-.055*\t)*(.45*cos(\ang)-.2*\t*sin(\ang)/sqrt(1+.04*\t*\t))},
         {.1*\t*\t+(.70-.055*\t)*(.22*cos(\ang)+sin(\ang)/sqrt(1+.04*\t*\t))});}
 \foreach \t in {0,5}{
  \draw[Ink,thick,domain=0:360,samples=80,variable=\a]
   plot ({\t+(.70-.055*\t)*(.45*cos(\a)-.2*\t*sin(\a)/sqrt(1+.04*\t*\t))},
         {.1*\t*\t+(.70-.055*\t)*(.22*cos(\a)+sin(\a)/sqrt(1+.04*\t*\t))});}
 \foreach \t/\lab in {0.7/1,2.5/2,4.2/3}{
  \draw[->,DeepBlue,very thick] (\t,{.1*\t*\t})--++(.55,{.11*\t});
  \draw[TextGray,thin] (\t,{.1*\t*\t})--(\t,{.1*\t*\t-1.0});
  \node[below,font=\small] at (\t,{.1*\t*\t-1.07}) {$\vect v_{\lab}$};}
 \node[left] at (-.5,0) {$C_1$};\node[left] at (-.5,-.45) {$S_A$};
 \node[right] at (5.55,2.8) {$C_2$};\node[right] at (5.55,2.35) {$S_B$};
\end{tikzpicture}\figcap{闭合截线 $C_1,C_2$ 围出横截面，侧壁由流线构成。箭头沿局部速度方向；图为弯曲流管的透视示意。}\label{ch13:fig:k10-2-24}\end{minipage}\end{center}}

\newcommand{\PipeContinuityDiagram}{%
\begin{center}\begin{minipage}{\linewidth}\centering\begin{tikzpicture}[scale=.74,>=Latex]
 \draw[very thick] (-4,.55)--(-1.5,.55)--(1.2,1.2)--(4,1.2);
 \draw[very thick] (-4,-.55)--(-1.5,-.55)--(1.2,-1.2)--(4,-1.2);
 \draw[vector] (-2.9,0)--(-1.65,0);
 \draw[TextGray,thin] (-2.3,-.12)--(-2.3,-.95);
 \node[below] at (-2.3,-1.0) {$v_A$};
 \draw[vector] (1.65,0)--(2.4,0) node[midway,below=4pt]{$v_B$};
 \draw[DeepBlue] (-3.25,-.55)--(-3.25,.55) node[above=5pt]{$S_A$};
 \draw[DeepBlue] (3.0,-1.2)--(3.0,1.2) node[above=5pt]{$S_B$};
\end{tikzpicture}\figcap{管道的两个横截面及流速方向，$v_B$ 为待求量。}\label{ch13:fig:k10-ex-2-97}\end{minipage}\end{center}}

\newcommand{\UniformPipeDiagram}{%
\begin{center}\begin{minipage}{\linewidth}\centering\begin{tikzpicture}[scale=.8,>=Latex]
 \draw[thick] (-4,.65)--(4,.65) (-4,-.65)--(4,-.65);
 \foreach \xx/\lab in {-2.8/1,2.0/2}{
  \draw[dashed,TextGray] (\xx,-.65)--(\xx,.65) node[above=5pt]{$S$};
  \draw[->,DeepBlue,thick] (\xx+.2,0)--(\xx+1.25,0) node[midway,below=5pt]{$v_{\lab}$};}
\end{tikzpicture}\figcap{同一圆柱管的两个等截面；在一维稳态模型中比较 $Sv_1$ 与 $Sv_2$。}\label{ch13:fig:k10-ex-2-98}\end{minipage}\end{center}}

\newcommand{\CompressiblePipeDiagram}{%
\begin{center}\begin{minipage}{\linewidth}\centering\begin{tikzpicture}[scale=.8,>=Latex]
 \draw[thick] (-4,.6)--(-1.5,.6)--(1.3,1.3)--(4,1.3);
 \draw[thick] (-4,-.6)--(-1.5,-.6)--(1.3,-1.3)--(4,-1.3);
 \foreach \xx/\rr/\lab in {-3/.6/A,3/1.3/B}{
  \draw[DeepBlue,dashed] (\xx,-\rr)--(\xx,\rr) node[above=6pt]{$S_{\lab}$};
  \node[below=6pt] at (\xx,-\rr) {$\rho_{\lab}$};}
 \draw[->,DeepBlue,thick] (-2.65,0)--(-1.65,0);
 \draw[TextGray,thin] (-1.9,.1)--(-1.9,1.05);
 \node[above] at (-1.9,1.1) {$v_A$};
 \draw[->,DeepBlue,thick] (1.55,0)--(2.5,0) node[midway,below=5pt]{$v_B$};
\end{tikzpicture}\figcap{两截面的面积、密度及流速符号，$v_B$ 为待求量。}\label{ch13:fig:k10-ex-2-100}\end{minipage}\end{center}}

\newcommand{\FallingJetDiagram}{%
\begin{center}\begin{minipage}{\linewidth}\centering\begin{tikzpicture}[scale=.85,>=Latex]
 \draw[thick] (-1.8,1)--(-1.8,0)--(-.50,0) (.50,0)--(1.8,0)--(1.8,1);
 \foreach \sgn in {-1,1}{\draw[DeepBlue,thick,domain=0:3.7,samples=90,variable=\s]
  plot ({\sgn*.5/(1+1.2*\s)^.25},{-\s});}
 \draw[TextGray,dashed] (-.75,0)--(1.3,0) (-.75,-3)--(1.3,-3);
 \draw[<->,TextGray] (-1.1,0)--(-1.1,-3) node[midway,left]{$h$};
 \draw[DeepBlue,dashed] (-.5,0)--(.5,0);
 \draw[DeepBlue,thick] ({-.5/4.6^.25},-3)--({.5/4.6^.25},-3);
 \node[above] at (0,.12) {$S$};\node[left] at (-.55,-3.25) {$S(h)$};
 \draw[->,DeepBlue,thick] (1.3,0)--(1.3,-.65) node[right]{$v$};
 \draw[->,DeepBlue,thick] (1.3,-3)--(1.3,-4.1) node[right]{$v(h)$};
\end{tikzpicture}\figcap{高度差 $h$ 的末端就是截面 $S(h)$；两速度箭头分别从对应截面起画。}\label{ch13:fig:k10-ex-2-99}\end{minipage}\end{center}}


\label{ch13:sec:flow-intro}

水和空气没有固定形状，受力后可以流动，统称流体。流体运动极其复杂，本节先讨论稳定流动：空间中任一给定点的速度大小和方向都不随时间改变。稳定并不表示某个流体微元的速度沿途不变；微元运动到不同位置时，会取得当地不同的流速。

流体由大量分子组成，逐一跟踪每个分子的运动通常既不可行也没有必要。稳定流动改为考察“空间中某一点此刻允许多大的速度、指向哪里”：只要这个速度场不随时间变化，后来经过同一点的流体微元就会取得相同速度。某个微元沿途经过不同位置时仍可能不断加速、减速或转弯，所以“稳定”描述的是空间分布不随时间变，并不等于每个微元作匀速直线运动。

水在通常条件下密度变化很小，可理想化为不可压缩流体；空气的密度较易随压强和速度改变，是可压缩流体。对稳定流动，给出空间每点的速度矢量便给出流速分布。沿速度场逐步跟踪微元可得到流线。

不可压缩同样是一种近似：它表示同一流体微元在运动中密度近似不变，并不是说流体“不能改变形状”。气体通常更容易压缩，完整描述稳定气流时除了速度场，还须给出各点密度。后面的连续性方程会分别说明体积流率和质量流率在两种模型中的区别。

\SphereStreamlinesDiagram

在速度场连续光滑且速度非零的区域，流线具有三条基本性质：流线在每一点与当地速度相切；每一点通过一条局部唯一的流线；两条不同流线不能相交。停滞点的速度为零，不能由零矢量指定切线方向，需作为例外单独讨论。

流线可以用逐步推进的方法构造：从一个起点出发，沿当地速度方向前进很短距离，到达新位置后改沿新位置的速度方向继续前进；令每一步趋于无限小，得到的光滑曲线就是流线。图~\ref{ch13:fig:k10-2-23}中的每条线都应在圆球附近平滑转弯而不穿过球体。若两条流线相交，交点处会出现两个切线方向，与稳定速度场在每点只有一个速度矢量矛盾。

取一个与流动方向垂直的面积 $S$，若截面上流速均为 $v$，单位时间通过的体积为 $Sv$。由一圈流线围成的管状区域称为流线管，其侧壁没有流体穿过。

在短时间 $\Delta t$ 内，穿过截面的流体近似占据底面积 $S$、长度 $v\Delta t$ 的柱体，体积为 $Sv\Delta t$，因此体积流率为 $Sv$。流线管的侧壁处处由流线组成，速度与侧壁相切，没有法向分量，所以流体不会从侧壁进出；流入和流出只能通过两端截面。

\StreamtubeDiagram

图~\ref{ch13:fig:k10-2-24}的两截面属于同一流线管。不可压缩流体稳定流过同一流线管时，单位时间通过各截面的体积相等。截面积增大则流速降低，截面积减小则流速增大。

\Needspace{16\baselineskip}\begin{origexamplebox}{不可压缩流体的连续性方程}\label{ch13:ex:k10-2-97}
不可压缩流体在管中稳定流动，每个垂直截面上流速均匀。截面 $A$ 的面积为 $S_A$、流速为 $v_A$；截面 $B$ 的面积为 $S_B$。求 $v_B$。

\PipeContinuityDiagram
\end{origexamplebox}
\Needspace{6\baselineskip}\begin{xesolution}
\ThAFlowVolume

图~\ref{ch13:fig:k10-ex-2-97}标出截面 $S_A,S_B$ 与对应速度。取相同时间 $\Delta t$，通过两截面的体积分别为 $S_Av_A\Delta t$、$S_Bv_B\Delta t$。稳态且不可压缩时两者相等：
\[
S_Av_A=S_Bv_B.
\]
所以
\[
\boxed{v_B=\frac{S_A}{S_B}v_A}.
\]
量纲检查：面积比无量纲，结果仍为速度。
\end{xesolution}

\Needspace{16\baselineskip}\begin{origexamplebox}{等截面管内能否稳定加速}\label{ch13:ex:k10-2-98}
圆柱形管道处处充满不可压缩流体。流体能否在管中作稳定的加速运动？采用沿直管轴线流动、各横截面流速均匀、无支流的一维模型。

\UniformPipeDiagram
\end{origexamplebox}
\Needspace{6\baselineskip}\begin{xesolution}
\ThAUniformFlow

图~\ref{ch13:fig:k10-ex-2-98}中圆柱管的截面积 $S$ 处处相同。连续性方程要求任意两截面
\[
Sv_1=Sv_2,
\]
故 $v_1=v_2$。在“不可压缩、稳定、处处充满且无支流”的条件下，流速不能沿管长增加，因而
\[
\boxed{\text{不能作沿程稳定加速运动}}.
\]
若流体可压缩、管道有泄漏或流动随时间变化，结论可能改变；这些都超出了题设。
若不作一维近似，连续性只保证截面平均流速相同；截面内部仍可发生速度重新分布。因此不能由“等截面”一句话推出任何微元的局部加速度都为零。
\end{xesolution}

\Needspace{16\baselineskip}\begin{origexamplebox}{下落水束的截面积}\label{ch13:ex:k10-2-99}
极大容器底部有面积为 $S$ 的圆孔，水刚离孔时速度为 $v$。每个水分子随后自由落体，水束不散开。求孔下方 $h$ 处水流截面积 $S(h)$。

\FallingJetDiagram
\end{origexamplebox}
\Needspace{6\baselineskip}\begin{xesolution}
\ThAFallingFlow

图~\ref{ch13:fig:k10-ex-2-99}中，$h$ 的下端、$S(h)$ 与 $v(h)$ 标在同一横截面。自由下落的速度由
\[
v(h)^2=v^2+2gh
\]
给出，故 $v(h)=\sqrt{v^2+2gh}$。水不可压缩，体积流率不变：
\[
Sv=S(h)v(h).
\]
因此
\[
\boxed{S(h)=\frac{Sv}{\sqrt{v^2+2gh}}}.
\]
随 $h$ 增大，水束加速而截面收细，乘积 $S(h)v(h)$ 保持为 $Sv$。
\end{xesolution}

对于可压缩流体，稳定流动中真正守恒的是质量流率而非体积流率。若没有物质产生或消失，对任意闭合曲面，单位时间流入的质量必须等于流出的质量，否则曲面内的质量会随时间积累或减少。这就是连续性条件。

若截面上密度和速度近似均匀，短时间通过的质量为 $\rho Sv\Delta t$，质量流率为 $\rho Sv$。不可压缩情形中同一流线管内 $\rho$ 不变，才可约去密度得到 $Sv$ 恒定；可压缩流体则可能用密度变化补偿截面积或速度变化。连续性只表达物质守恒，不自动给出压强和速度的全部关系，解决更完整的流体问题还需要后续动力学方程。

连续性方程的适用对象必须选成同一股稳定流。若管道存在支路，应把所有入口和出口的质量流率带符号相加，而不能只随意比较两个截面；若流动随时间变化，控制体内还可能发生质量积累。对题图的单一流线管，侧壁无穿流、内部又无源汇，所以任意两个横截面的流率才相等。

自由水束的两个例题展示了同一守恒条件的两种实现方式。真实水近似不可压缩，密度保持不变，下落加速后必须缩小截面积；若人为规定截面积不变而又让各微元自由加速，连续性就要求密度降低。二者不能同时把截面积和密度都固定，否则 $v$ 增大将使质量流率凭空增大，违反稳定流动的物质守恒。

定常流动中，流线与单个微元的轨迹重合；非定常流动中一般不重合。若速度场随时间改变，某一时刻画出的流线仍处处与当时速度相切，但一个微元随后看到的速度场已经变化，它实际走出的迹线可能与这条瞬时流线不同。本节始终假定稳定流动，因此可以用流线同时表示当地速度方向和微元运动轨迹；超出这一条件时应区分两者。

截面若不垂直于速度，穿过它的体积流率应使用法向分量 $v_n$，即 $Sv_n$，而不是总速率 $Sv$。若截面上速度或密度不均匀，还需把 $\rho\vect v\cdot\vect n$ 在整个截面上积分。本节例题把截面选为垂直于流动且各处均匀，才得到简洁的 $Sv$ 或 $\rho Sv$；这些条件也是公式的一部分，不能从图中省略。

\Needspace{16\baselineskip}\begin{origexamplebox}{可压缩流体的质量连续性}\label{ch13:ex:k10-2-100}
流体在管道中稳定流动，且同一垂直截面上性质均匀。$A$ 处截面积、密度、流速分别为 $S_A,\rho_A,v_A$；$B$ 处为 $S_B,\rho_B,v_B$。求 $v_B$。

\CompressiblePipeDiagram
\end{origexamplebox}
\Needspace{6\baselineskip}\begin{xesolution}
\ThAFlowMass

图~\ref{ch13:fig:k10-ex-2-100}将两侧密度分别标为 $\rho_A,\rho_B$。单位时间通过截面的质量为 $\rho Sv$，不能只比较体积。质量连续性给出
\[
\rho_AS_Av_A=\rho_BS_Bv_B,
\]
所以
\[
\boxed{v_B=\frac{\rho_AS_A}{\rho_BS_B}v_A}.
\]
当 $\rho_A=\rho_B$ 时退化为例~\ref{ch13:ex:k10-2-97}的不可压缩结果。
\end{xesolution}

\Needspace{16\baselineskip}\begin{origexamplebox}{等截面自由落体液柱的密度}\label{ch13:ex:k10-2-101}
极大容器中液体初始密度为 $\rho$，从面积为 $S$ 的孔以速度 $v$ 流出。忽略分子间相互作用和表面张力，设下落过程中水束截面积始终为 $S$，每个分子作自由落体。求密度 $\rho(h)$。

\ThQDilutingColumn
\end{origexamplebox}
\Needspace{6\baselineskip}\begin{xesolution}
\ThADilutingColumn

下落 $h$ 后速度仍为
\[
v(h)=\sqrt{v^2+2gh}.
\]
稳态质量流率在孔口和高度 $h$ 处相等：
\[
\rho Sv=\rho(h)S v(h).
\]
截面积约去，得到
\[
\boxed{\rho(h)=\frac{\rho v}{\sqrt{v^2+2gh}}}.
\]
这里人为固定了截面积，因此加速必须由密度降低来满足质量连续性；这正与例~\ref{ch13:ex:k10-2-99}中“密度不变、截面积减小”互为对照。
\end{xesolution}
}

\newcommand{\FluidBernoulliDiagram}{%
\begin{center}\begin{minipage}{\linewidth}\centering\begin{tikzpicture}[scale=.9,>=Latex]
 \draw[->,TextGray] (-2,-2.6)--(-2,3.4) node[above]{$z$};
 \draw[TextGray,dashed] (-2,-2.3)--(7.8,-2.3);
 \begin{scope}[rotate=16]
  \draw[thick] (0,.85)--(6,.5) (0,-.85)--(6,-.5);
  \draw[TextGray,dashed] (0,0)--(6,0);
  \draw[DeepBlue,thick] (0,0) ellipse[x radius=.20,y radius=.85];
  \draw[DeepBlue,thick] (6,0) ellipse[x radius=.117647,y radius=.5];
  \draw[TextGray,dashed] (.35,0) ellipse[x radius=.195294,y radius=.830];
  \draw[TextGray,dashed] (7.0115,0) ellipse[x radius=.117647,y radius=.5];
  \draw[->,DeepBlue,thick] (-1.1,0)--(0,0);
  \draw[->,DeepBlue,thick] (7.4,0)--(6,0);
  \node[above] at (-.9,.15) {$p_1A_1$};\node[above right] at (7.0,.6) {$p_2A_2$};
  \draw[TextGray,thin] (0,-.87)--(0,-1.28) (.35,-.87)--(.35,-1.28);
  \draw[->,TextGray] (-.55,-1.15)--(0,-1.15);
  \draw[->,TextGray] (.90,-1.15)--(.35,-1.15);
  \node[below] at (.175,-1.3) {$\dd s_1$};
  \draw[<->,TextGray] (6,-.85)--(7.0115,-.85);
  \node[below] at (6.5,-1.0) {$\dd s_2$};
  \node[above] at (0,1.02) {$A_1,\,v_1$};\node[above] at (5.9,.72) {$A_2,\,v_2$};
  \coordinate (FluidA) at (0,0);\coordinate (FluidB) at (6,0);
 \end{scope}
 \draw[<->,TextGray] (2,-2.3)--(2,0) node[midway,right]{$z_1$};
 \draw[TextGray,dotted] (FluidA)--(2,0);
 \draw[<->,TextGray] (5,-2.3)--(5,{6*sin(16)}) node[midway,left]{$z_2$};
 \draw[TextGray,dotted] (FluidB)--(5,{6*sin(16)});
\end{tikzpicture}\figcap{流管两端压力做功。相同时间内移入、移出的质量相等；高度从同一水平基准面量起。}\label{ch13:fig:bernoulli}
\end{minipage}\end{center}}

\newcommand{\FluidShearDiagram}{%
\begin{center}\begin{minipage}{\linewidth}\centering\begin{tikzpicture}[scale=.85,>=Latex]
 \foreach \yy in {0,1.8}{\draw[Ink,thick] (0,\yy)--(4.2,\yy)--(5.15,\yy+.55)--(.95,\yy+.55)--cycle;}
 \foreach \yy in {.15,.60,1.05,1.50}{\draw[->,DeepBlue] (1,\yy)--({1.65+.65*\yy},\yy);}
 \draw[->,DeepBlue,very thick] (2,2.08)--(4.25,2.08);
 \node[right] at (5.5,2.08) {$v_x+\dd v_x$};
 \draw[->,DeepBlue,very thick] (2,.27)--(3.4,.27) node[right=4pt]{$v_x$};
 \draw[<->,TextGray] (-.45,0)--(-.45,1.8) node[midway,left]{$\dd y$};
 \draw[->,Ink] (-1.3,0)--(-1.3,2.7) node[above]{$y$};
 \draw[->,Ink] (0,-.8)--(2,-.8) node[right]{$x$};
 \node[above] at (1.6,2.5) {上层对下层：向右拉};
 \node[below] at (3.1,-1.05) {$\tau_{xy}=\dd F_x/\dd S=\eta\,\dd v_x/\dd y$};
\end{tikzpicture}\figcap{相距 $\dd y$ 的流层与切向速度梯度。上层流得较快，黏性作用减小两层速度差；反作用力方向相反。}\label{ch13:fig:shear}
\end{minipage}\end{center}}

\newcommand{\FluidOutletDiagram}{%
\begin{center}\begin{minipage}{\linewidth}\centering\begin{tikzpicture}[scale=.78,>=Latex]
 \fill[Blue!12] (0,0) rectangle (4,3);
 \draw[thick] (0,3.6)--(0,0)--(4,0)--(4,.7) (4,1.05)--(4,3.6);
 \draw[Blue,thick] (0,3)--(4,3);
 \draw[Blue,thick] (4,.7)--(5.5,.5) (4,1.05)--(5.5,.85);
 \fill (1,3) circle(1.6pt) node[above]{$A$};
 \fill (4,.875) circle(1.6pt);
 \draw[TextGray,thin] (4,.875)--(3.4,1.3);
 \node[above left] at (3.4,1.3) {$B$};
 \draw[->,DeepBlue,thick] (4.7,.67)--(5.6,.67) node[right]{$v$};
 \draw[<->,TextGray] (-.6,.875)--(-.6,3) node[midway,left]{$h$};
 \draw[dotted,TextGray] (-.6,.875)--(3.7,.875);
 \node[above] at (2,4.0) {$p_A=p_B=p_0$};
\end{tikzpicture}\figcap{取自由液面 $A$ 与小孔出口 $B$ 比较压强、速度和高度。}\label{ch13:fig:outlet}\end{minipage}\end{center}}

\newcommand{\FluidPlateDiagram}{%
\begin{center}\begin{minipage}{\linewidth}\centering\begin{tikzpicture}[scale=.85,>=Latex]
 \draw[Ink,very thick] (4,-2)--(4,2);
 \draw[Blue,thick] (0,-.25)--(3.5,-.25)..controls(3.8,-.25)and(3.8,-1)..(3.8,-1.65);
 \draw[Blue,thick] (0,.25)--(3.5,.25)..controls(3.8,.25)and(3.8,1)..(3.8,1.65);
 \draw[->,DeepBlue,thick] (.5,0)--(2,0) node[midway,above=8pt]{$v$};
 \draw[->,DeepBlue,thick] (3.65,.7)--(3.65,1.65);
 \draw[->,DeepBlue,thick] (3.65,-.7)--(3.65,-1.65);
 \draw[dashed,TextGray] (2.3,-1.75) rectangle (4.3,1.75);
 \draw[->,Ink,thick] (4.7,0)--(6,0) node[right]{$F_{\text{水对板}}$};
 \node[left] at (0,0) {$A$};\node[above] at (3.3,2.2) {虚线为控制体};
 \draw[->,TextGray] (0,-1.6)--(1.4,-1.6) node[right]{$x$};
\end{tikzpicture}\figcap{来流有正向 $x$ 动量，流出水沿板面，法向动量为零。}\label{ch13:fig:jet-plate}\end{minipage}\end{center}}

\newcommand{\FluidSettlingDiagram}{%
\begin{center}\begin{minipage}{\linewidth}\centering\begin{tikzpicture}[scale=.9,>=Latex]
 \draw[thick] (0,0) circle(.62);
 \draw[->,Ink,thick] (0,-.62)--(0,-2) node[right]{$mg$};
 \draw[->,DeepBlue,thick] (-.26,{sqrt(.62*.62-.26*.26)})--(-.26,1.6) node[above left]{$F_{\!B}$};
 \draw[->,DeepBlue,thick] (.26,{sqrt(.62*.62-.26*.26)})--(.26,2.2) node[above right]{$6\pi\eta Rv_t$};
 \draw[->,TextGray,thick] (1.8,.6)--(1.8,-.7) node[right]{$v_t$};
 \node[right] at (2.2,1.6) {$F_{\!B}=\frac43\pi R^3\rho_f g$};
\end{tikzpicture}\figcap{下沉小球的受力。终端状态速度向下，浮力和黏性阻力均向上。}\label{ch13:fig:settling}\end{minipage}\end{center}}

\section{流体运动学与物质导数}\label{ch13:sec:material}

流体由大量不断更换邻居的质元组成。拉格朗日描述（Lagrangian description）给每个质元贴上初始标签 $\vect a$，记录其位置 $\vect r(\vect a,t)$；欧拉描述（Eulerian description）则在空间各点记录速度、密度和压强场
\begin{equation}
 \vect v=\vect v(\vect r,t),\qquad \rho=\rho(\vect r,t),\qquad p=p(\vect r,t).
\end{equation}
两种描述说的是同一群质元。固定初始标签 $\vect a$，对其位置求时间导数，就得到该质元此刻的速度：
\[
 \vect v\bigl(\vect r(\vect a,t),t\bigr)
 =\left.\frac{\partial\vect r(\vect a,t)}{\partial t}\right|_{\vect a}.
\]
左边先找到质元当前所在的位置，再读取空间速度场；右边始终跟踪同一个质元。标签 $\vect a$ 不能在求导时随意改变。
竞赛流体题通常使用欧拉描述，但研究某个质元所经历的变化时必须沿它的轨迹求导。若标量场为 $f(\vect r,t)$，物质导数（material derivative）为
\begin{equation}
 \frac{\mathrm D f}{\mathrm Dt}
 =\frac{\partial f}{\partial t}+\vect v\cdot\nabla f.
\end{equation}
这是多元函数链式法则的直接结果。把沿质元轨迹读到的量写成
$f(x(t),y(t),z(t),t)$，则
\[
 \frac{\dd f}{\dd t}=\frac{\partial f}{\partial t}
 +\frac{\partial f}{\partial x}\frac{\dd x}{\dd t}
 +\frac{\partial f}{\partial y}\frac{\dd y}{\dd t}
 +\frac{\partial f}{\partial z}\frac{\dd z}{\dd t}.
\]
代入 $\dd x/\dd t=v_x$ 等关系，后三项才合写成 $\vect v\cdot\nabla f$。
第一项描述固定位置随时间的变化，后三项描述质元移动到不同位置而读到的变化；二者都不能遗漏。
因此流体质元的加速度不是单纯的 $\partial\vect v/\partial t$，而是
\begin{equation}
 \vect a=\frac{\mathrm D\vect v}{\mathrm Dt}
 =\frac{\partial\vect v}{\partial t}+(\vect v\cdot\nabla)\vect v.
\end{equation}
第二项是对流加速度：即使某点的速度场不随时间改变，质元进入速度不同的区域时仍可有加速度。

迹线是一个质元随时间走过的轨迹，满足 $\dd\vect r/\dd t=\vect v(\vect r,t)$；流线是在同一时刻处处与速度相切的曲线。定常流动中二者重合；非定常流动中一般不重合，非定常流动时不能把某一时刻画出的流线当成质元未来轨迹。
对固定时刻 $t_0$，流线可用参数 $s$ 表示为 $\dd\vect r/\dd s=\lambda(s)\vect v(\vect r,t_0)$。
在相应分量非零处，也常写成
\[
 \frac{\dd x}{v_x(x,y,z,t_0)}=\frac{\dd y}{v_y(x,y,z,t_0)}
 =\frac{\dd z}{v_z(x,y,z,t_0)}.
\]
某分量为零时应使用参数形式，不能除以零。前面的流管图中，$C_1,C_2$ 之间的侧壁由同一组流线组成，速度沿侧壁而不穿过它。

为比较相邻质元，令两位置相差 $\delta\vect r$，作一阶展开：
\[
 v_i(\vect r+\delta\vect r)=v_i(\vect r)
 +\sum_j\frac{\partial v_i}{\partial x_j}\delta x_j+O(|\delta\vect r|^2).
\]
第一项给出共同平动。剩余的速度梯度可分解成对称和反对称两部分：
\begin{equation}
 \frac{\partial v_i}{\partial x_j}
 =D_{ij}+\Omega_{ij},
 \quad
 D_{ij}=\frac12\left(\frac{\partial v_i}{\partial x_j}+\frac{\partial v_j}{\partial x_i}\right),
 \quad
 \Omega_{ij}=\frac12\left(\frac{\partial v_i}{\partial x_j}-\frac{\partial v_j}{\partial x_i}\right).
\end{equation}
$D_{ij}$ 描述局部拉伸和剪切速率，$\Omega_{ij}$ 描述刚体式转动；涡量 $\vect\omega=\nabla\times\vect v$，局部流体微团的角速度为 $\vect\omega/2$。不可压缩条件是 $\nabla\cdot\vect v=\operatorname{tr}\mathbf D=0$，表示质元体积不随运动改变，并不等价于“速度处处相同”。
具体地，$\boldsymbol\Omega\,\delta\vect r=(\vect\omega/2)\times\delta\vect r$，所以它只改变相邻质元连线的方向。
连线长度平方的变化率则为
\[
 \frac{\dd}{\dd t}|\delta\vect r|^2
 =2\,\delta\vect r\cdot\mathbf D\,\delta\vect r,
\]
反对称部分的贡献为零。故 $\mathbf D$ 才描述变形快慢；它的量纲是 $\mathrm{s^{-1}}$，不同于弹性体中无量纲的应变。

\section{静止流体}\label{ch13:sec:statics}

静止流体不能承受持续剪切应力：一旦持续受切向力，普通流体会不断变形，而不是保持静止。
任取内部一个面，其表面力只能沿法线；如果不同方向的正应力不等，换一个倾斜截面便会出现切向分量，与静止矛盾。
所以静止流体的应力各向同性。取拉伸应力为正，应力张量为
\[
 \boldsymbol\sigma=-p\mathbf I.
\]
体积 $V$ 内流体所受周围流体的压力合力为
\[
 \vect F_p=-\oint_{\partial V}p\vect n\,\dd S
 =-\int_V\nabla p\,\dd V,
\]
其中 $\vect n$ 是向外单位法向，第二步对各分量使用高斯定理。因此单位体积的压力合力为 $-\nabla p$。
若外体积力密度为 $\vect f$，静力平衡要求 $\nabla p=\vect f$。这还意味着在光滑单连通区域内
$\nabla\times\vect f=0$；注意 $\vect f$ 是单位体积的力，不能与单位质量外力混淆。
在重力场中 $\vect f=\rho\vect g$，取竖直向上为 $z$，薄层受力平衡给出
\begin{equation}
 \frac{\dd p}{\dd z}=-\rho g.
\end{equation}
不可压缩液体密度恒定时，深度增加 $h$，压强增加 $\rho gh$。同一连通静止液体中，同一水平面压强相等；
平直界面或忽略表面张力时，不同液体交界面两侧压强连续；曲面界面还可能存在后文讨论的拉普拉斯压差。
物体浸入液体后，液体对其表面的压力合力等于把相同体积液体移走所失去的压力合力：
\[
 \vect F_B=-\int_{V_{\text{排}}}\rho\vect g\,\dd V.
\]
故浮力等于排开流体的重量、方向与重力相反。密度均匀时，其作用线通过排开体积的几何中心；密度不均匀时应取排开流体的质量中心。

\subsection{从一小柱液体看压强随高度的变化}
压强不是向下的力，而是作用在所选小面上的法向力除以面积。液体内部同一点的压强是标量，受压面的朝向不同，压力方向随之改变。因此不能因为重力向下，就把侧壁受到的液体压力也画成向下。

取截面积为 $S$、底面高度为 $z$、顶面高度为 $z+\dd z$ 的竖直小柱。周围液体在底面给它向上的力 $p(z)S$，在顶面给它向下的力 $p(z+\dd z)S$；侧面压力没有竖直分量。向上为正，平衡式为
\[
 p(z)S-p(z+\dd z)S-\rho gS\dd z=0.
\]
减去相同项、除以 $S\dd z$ 并取极限，才得到 $\dd p/\dd z=-\rho g$。负号表示向上压强降低；若改用向下的深度 $h$，则 $\dd p/\dd h=\rho g$。对密度恒定的液体，从自由面 $p=p_a$ 积分到深度 $h$，有
\[
 p(h)-p_a=\int_0^h\rho g\,\dd h'=\rho gh.
\]
$p$ 是绝对压强，$p-p_a$ 是相对大气的表压。自由面与大气相通并不意味着 $p=0$，而是表压为零。比较受压面两侧的合力时，才可能把相同的大气压贡献抵消。

\subsection{不同液层与受压面的合力}
密度分层时，应在每层分别累加压强增量。在同一连通静止液体中取水平小柱，水平外体积力为零，压强差必须为零；这就是等高处等压的力学依据。若两点属于互不连通的液体，或者有水平加速度，就不能仅凭高度相同认定压强相等。

\Needspace{16\baselineskip}\begin{origexamplebox}{两层液体中的侧壁压力}\label{ch13:ex:teaching-layered}
一个敞口容器中，上层液体密度为 $\rho_1$、厚度为 $h_1$，下层密度为 $\rho_2>\rho_1$、厚度为 $h_2$。竖直侧壁宽度为 $b$，外侧为大气。求液体作用在侧壁下层部分上的净压力及其作用点深度，忽略界面张力。

\ThQLayers
\end{origexamplebox}
\Needspace{6\baselineskip}\begin{xesolution}
\ThALayers

从两液体界面向下量 $x$，$0\le x\le h_2$。上层已贡献表压 $\rho_1gh_1$，下层再贡献 $\rho_2gx$，故
\[
 \Delta p(x)=\rho_1gh_1+\rho_2gx,
 \qquad \dd F=\Delta p(x)b\,\dd x.
\]
每一水平窄条所受净压力均沿侧壁法线向外，因而可以直接相加其大小：
\[
 F=bg\left(\rho_1h_1h_2+\frac12\rho_2h_2^2\right).
\]
求作用点不能直接取几何中心，因为下方压强更大。设合力位于界面下方 $x_c$，对界面处水平轴取矩，有
\[
 Fx_c=\int_0^{h_2}x\,\dd F
 =bg\left(\frac12\rho_1h_1h_2^2+\frac13\rho_2h_2^3\right),
\]
所以
\[
 x_c=\frac{\rho_1h_1h_2^2/2+\rho_2h_2^3/3}
 {\rho_1h_1h_2+\rho_2h_2^2/2}.
\]
距自由面的深度是 $h_1+x_c$。没有上层时 $x_c=2h_2/3$；若上层压力远大于下层压强变化，分布接近均匀，$x_c\to h_2/2$。两极限分别对应三角形和近似矩形压强分布，且 $\rho gbh^2$ 的量纲确为力。
\end{xesolution}

\subsection{浮力来自压力差，稳定性还取决于力矩}
设想把浸没物体占据的空间重新填满周围液体。静止时，这部分液体的重量被周围压力合力平衡；只要边界上的压强分布相同，把它换成物体后，压力合力也不变。这一替换论证给出浮力，而不是另外假设一种向上的力。
对于跨越两层液体的物体，应分别计算 $\rho_1gV_1+\rho_2gV_2$，不能把总排水体积都乘同一个密度。浮力与重力相等只保证竖直合力为零；物体稍微倾斜后是否产生回复力矩，还需比较重心与浮力作用线，不能仅据“浮着”判断稳定。

\Needspace{16\baselineskip}\begin{origexamplebox}{浮体的竖直小振动}\label{ch13:ex:teaching-buoyancy}
质量为 $m$、水平截面积恒为 $S$ 的直立柱体，在密度 $\rho$ 的液体中平衡漂浮。只允许竖直运动，忽略流体附加惯性、阻力和表面张力，且小振动中底面始终浸没、顶面始终露出。求平衡浸深及小振动周期。

\ThQFloat
\end{origexamplebox}
\Needspace{6\baselineskip}\begin{xesolution}
\ThAFloat

平衡时 $mg=\rho gSh_0$，所以 $h_0=m/(\rho S)$。取柱体向下的位移 $x$ 为正；此时浸深为 $h_0+x$，浮力向上且增大：
\[
 m\ddot x=mg-\rho gS(h_0+x)=-\rho gSx.
\]
平衡项消去后，剩下的正是指向平衡位置的回复力。因此
\[
 \omega^2=\frac{\rho gS}{m},\qquad
 T=2\pi\sqrt{\frac{m}{\rho gS}}.
\]
若 $x<0$，浮力减少，合力向下，也把柱体推回平衡。$\rho gS$ 的单位为 $\mathrm{N/m}$，相当于等效劲度系数。真实液体随物体一起加速会增加有效惯性，故本题结果依赖题设的忽略条件；也不能把它用于已经完全浸没、排液体积不随 $x$ 变化的柱体。
\end{xesolution}

\section{连续性方程与伯努利方程}\label{ch13:sec:bernoulli}

一般质量守恒写成
\begin{equation}
 \frac{\partial\rho}{\partial t}+\nabla\cdot(\rho\vect v)=0.
\end{equation}
对固定控制体，质量的增加率等于净流入率：
\[
 \frac{\dd}{\dd t}\int_V\rho\,\dd V
 =-\oint_{\partial V}\rho\vect v\cdot\vect n\,\dd S.
\]
把右边化成体积分，利用控制体可以任意选取，才得到上面的局部式。
再展开 $\nabla\cdot(\rho\vect v)=\vect v\cdot\nabla\rho+\rho\nabla\cdot\vect v$，可写成
\[
 \boxed{\frac{\mathrm D\rho}{\mathrm Dt}+\rho\nabla\cdot\vect v=0}.
\]
这说明散度正负控制质元体积的胀缩；质量不变时密度随之降低或升高。不可压缩流体满足
$\mathrm D\rho/\mathrm Dt=0$，因而 $\nabla\cdot\vect v=0$；不同质元原本可以有不同密度。
对稳定的一维流管，质量流率 $\dot m=\rho Av$ 沿程不变；不可压缩时化为 $Av=\text{常量}$。

无黏、不可压缩流体沿同一条流线作稳定流动时，欧拉方程（Euler equation）沿流线积分得到
\begin{equation}
 p+\frac12\rho v^2+\rho gz=\text{常量}.
\end{equation}
三项分别是单位体积的压强能、动能和重力势能。伯努利方程（Bernoulli equation）不能跨越泵、涡轮、强黏性耗散区随意使用，
也不能在旋涡流中默认所有流线具有同一常量。
下面把推导写开。对无黏质元应用牛顿定律，压力合力密度为 $-\nabla p$，故
\[
 \rho\frac{\mathrm D\vect v}{\mathrm Dt}=-\nabla p+\rho\vect g.
\]
在定常流中，沿流线取弧长 $s$，切向加速度为 $v\,\dd v/\dd s$。沿切向投影得到
\[
 \rho v\frac{\dd v}{\dd s}=-\frac{\dd p}{\dd s}-\rho g\frac{\dd z}{\dd s},
 \qquad \dd(v^2/2)+\frac{\dd p}{\rho}+g\,\dd z=0.
\]
密度为常数时逐项积分，就是伯努利方程。
\FluidBernoulliDiagram
图~\ref{ch13:fig:bernoulli}也给出功—能解释。相同短时间内，两端移过的质量相等：
$\dd m=\rho A_1\dd s_1=\rho A_2\dd s_2$。
两端压力总功为
\[
 \dd W_p=p_1A_1\dd s_1-p_2A_2\dd s_2
 =\dd m\left(\frac{p_1}{\rho}-\frac{p_2}{\rho}\right).
\]
它等于动能与重力势能的增加
$\dd m[(v_2^2-v_1^2)/2+g(z_2-z_1)]$。
出口压力做负功，正是式中压强项不能漏掉的原因。

更一般地，对定常、无黏、正压流体，若密度可写成压强的单值函数 $\rho=\rho(p)$，沿流线有
\begin{equation}
 \frac{v^2}{2}+gz+\int^p\frac{\dd p'}{\rho(p')}=\text{常量}.
\end{equation}
不可压缩时 $\rho$ 为常量，积分项化为 $p/\rho$；可压缩气流中密度和内能通常与温度耦合，不能仍把 $p+\rho v^2/2+\rho gz$ 当作普适守恒式。
若是定常、无黏、绝热且没有激波等不可逆过程的流动，单位质量内能 $u$ 满足
\[
 \dd u=-p\,\dd(1/\rho)=\frac{p}{\rho^2}\dd\rho.
\]
定义单位质量焓 $h_{\rm sp}=u+p/\rho$，则
\[
 \dd h_{\rm sp}=\dd u+\frac{\dd p}{\rho}-\frac p{\rho^2}\dd\rho
 =\frac{\dd p}{\rho}.
\]
代回沿流线的动力学式，得到
\[
 \boxed{h_{\rm sp}+\frac{v^2}{2}+gz=\text{沿该流线的常量}}.
\]
这里焓含有内能，不能只保留 $p/\rho$。例如理想气体的内能随温度改变，气流加速可能伴随温度降低。
受热驱动的自然对流还涉及密度变化和浮力；液体密度变化虽小，也不能一概删去其浮力效应。
常用近似是在质量守恒中取无散速度场，而在重力项中保留小密度差；这与“所有动力学都与温度无关”不同。

\Needspace{16\baselineskip}\begin{origexamplebox}{小孔流出速度}\label{ch13:ex:outlet}
大水箱自由液面到侧壁小孔的高度差为 $h$，水箱横截面远大于小孔。忽略黏性，求流出速度。

\ThQOutlet
\end{origexamplebox}
\Needspace{6\baselineskip}\begin{xesolution}
在自由液面与小孔出口之间使用伯努利方程。两点都与大气相通，压强项相消；大水箱自由液面速度近似为零：
\FluidOutletDiagram
图~\ref{ch13:fig:outlet}中 $z_A-z_B=h$，$v_A\simeq0$，于是
\[
 p_0+\rho gh=p_0+\frac12\rho v^2.
\]
所以
\[
 \boxed{v=\sqrt{2gh}}.
\]
这与物体自由下落高度 $h$ 的速度相同，但流体微元的真实轨迹和压强变化并不等同于一个孤立质点自由落体。
\end{xesolution}

\section{动量流与流体作用力}\label{ch13:sec:momentum}

稳定流动中，控制体所受外力等于流出动量流率减去流入动量流率：
\begin{equation}
 \sum\vect F=\sum_{\rm out}\dot m\vect v-
 \sum_{\rm in}\dot m\vect v.
\end{equation}
这里的外力包括压强、重力和固体壁面对流体的力。求流体对管壁的力时，最后还要用牛顿第三定律反向。

\Needspace{16\baselineskip}\begin{origexamplebox}{水射流正撞平板}\label{ch13:ex:jet-plate}
密度为 $\rho$、截面积为 $A$、速度为 $v$ 的稳定水射流垂直撞击静止大平板，撞后水沿板面向四周流走。忽略重力，求水对板的平均力。

\ThQJet
\end{origexamplebox}
\Needspace{6\baselineskip}\begin{xesolution}
取法向为 $x$ 轴。质量流率 $\dot m=\rho Av$。流入法向速度为 $v$，流出水流沿板面，法向速度为零。
\FluidPlateDiagram
图~\ref{ch13:fig:jet-plate}中虚线包围转向区；图右箭头是水对板的力，列控制体方程时应使用它的反作用力。
平板对水的力为
\[
 F_{\text{板对水},x}=\dot m(0-v)=-\rho Av^2.
\]
所以水对板的力方向沿来流方向，大小
\[
 \boxed{F=\rho Av^2}.
\]
量纲为 $\mathrm{kg\,m^{-3}}\cdot\mathrm{m^2}\cdot\mathrm{m^2s^{-2}}=\mathrm N$。
\end{xesolution}

\subsection{为什么稳态流动仍然需要外力}
稳态表示固定位置的流速不随时间变化，不表示经过这里的每一份流体都没有加速度。图~\ref{ch13:fig:jet-plate} 中，虚线区域内的水总量可以不变，但新进入的水不断丢失法向动量，这就需要平板持续施力。
在短时间 $\dd t$ 内，流入质量为 $\dot m\dd t$，流入动量为 $\dot m\vect v_{\rm in}\dd t$；流出端同理。稳态时控制体内的动量不积累，所以外力冲量必须补足出入动量之差。除以 $\dd t$ 得到上面的动量流方程。多出口时必须先按矢量相加，不能直接相加速率。

对管内流动，在入口截面，外部流体把所研究流体推入控制体；在出口截面，外部流体阻碍它流出。两个压力方向均指向控制体内部。若求管道与环境之间的净作用力，使用表压时还须把管外大气压的抵消计入，不能只把绝对压强随意置零。

\Needspace{16\baselineskip}\begin{origexamplebox}{射流推动匀速后退的平板}\label{ch13:ex:teaching-moving-plate}
沿水平 $+x$ 方向的细水射流速度为 $v$、截面积为 $A$、密度为 $\rho$。一块足够大的平板沿同方向以恒速 $u$ 后退，$0\le u<v$。水撞板后沿板面流走，相对平板没有法向速度。忽略重力及空气作用，求水对板的力和传给板的功率。

\ThQMovingPlate
\end{origexamplebox}
\Needspace{6\baselineskip}\begin{xesolution}
\ThAMovingPlate

以匀速平板为参考系，它仍是惯性系。对应图~\ref{ch13:fig:jet-plate} 的控制体随板平移，入射相对速度为 $v-u$，故每秒真正撞板的质量为
\[
 \dot m_{\rm hit}=\rho A(v-u).
\]
这不是固定截面的 $\rho Av$：一部分水还在追赶平板。水相对板的法向速度从 $v-u$ 变成零，每单位质量丢失法向动量 $v-u$。由反作用关系，板受到沿 $+x$ 的力
\[
 F=\dot m_{\rm hit}(v-u)=\rho A(v-u)^2.
\]
也可在地面系看同一份水，撞前法向速度为 $v$，撞后为 $u$，其动量变化给出同一个结果。力对平板的功率必须使用地面系板速：
\[
 P=Fu=\rho Au(v-u)^2.
\]
在 $0<u<v$ 内，$\dd P/\dd u=\rho A(v-u)(v-3u)$，故最大功率出现在 $u=v/3$，大小为 $4\rho Av^3/27$。$u=0$ 时有力却不向静止板传机械功；$u\to v$ 时几乎没有水追上平板，力和功率均趋零。这也说明力最大与功率最大不是同一个条件。
\end{xesolution}

\section{黏性、层流与相似准则}\label{ch13:sec:viscosity}

牛顿流体（Newtonian fluid）的剪切应力与速度梯度成正比：
\begin{equation}
 \tau=\eta\frac{\partial v_x}{\partial y},
\end{equation}
其中 $\eta$ 为动力黏度。半径 $R$、长度 $L$ 的圆管内充分发展层流满足哈根--泊肃叶定律（Hagen--Poiseuille law）
\begin{equation}
 Q=\frac{\pi R^4}{8\eta L}\Delta p.
\end{equation}
\FluidShearDiagram
图~\ref{ch13:fig:shear}中，上下两流层沿 $x$ 流动，速度只随法向坐标 $y$ 改变。
速度差除以层间距是剪切速率；乘以动力黏度才是单位面积的切向力。
$\eta$ 的单位为 $\mathrm{Pa\,s}$，另定义运动黏度
\[
 \boxed{\nu=\eta/\rho},\qquad [\nu]=\mathrm{m^2/s}.
\]
推导圆管流量时，设 $G=\Delta p/L>0$，距轴 $r$ 处的定常速度为 $v_z(r)$。沿轴的力平衡给出
$\eta r^{-1}\dd[r\,\dd v_z/\dd r]/\dd r=-G$。
积分一次为 $r\,\dd v_z/\dd r=-Gr^2/(2\eta)+C$；轴心速度梯度有限要求 $C=0$。
再积分并用管壁无滑移条件 $v_z(R)=0$：
\[
 v_z(r)=\frac{G}{4\eta}(R^2-r^2),\qquad
 Q=\int_0^R v_z(r)\,2\pi r\dd r
 =\frac{\pi GR^4}{8\eta}.
\]
这一步同时说明充分发展、直圆管、不可压缩牛顿流体和无滑移边界等条件为何重要。
流量对半径的四次方依赖说明细管尺寸误差会产生很大影响。雷诺数（Reynolds number）
\begin{equation}
 \mathrm{Re}=\frac{\rho v\ell}{\eta}
\end{equation}
比较惯性与黏性效应；其临界值依赖几何和扰动，不能把某个数值当成所有流动的普适分界。

低雷诺数下，半径 $R$ 的球以速度 $v$ 在黏性流体中缓慢运动，斯托克斯阻力（Stokes drag）为 $6\pi\eta Rv$。

\Needspace{16\baselineskip}\begin{origexamplebox}{小球在黏性液体中的终端速度}\label{ch13:ex:settling}
密度为 $\rho_s$、半径为 $R$ 的小球在密度 $\rho_f$、黏度 $\eta$ 的液体中下落。假设满足斯托克斯阻力，求终端速度。

\ThQSettling
\end{origexamplebox}
\Needspace{6\baselineskip}\begin{xesolution}
终端状态合力为零。向下重力减去浮力等于向上黏性阻力：
\FluidSettlingDiagram
图~\ref{ch13:fig:settling}按下沉情形画出三个力；向下为正，平衡式为
\[
 \frac43\pi R^3(\rho_s-\rho_f)g=6\pi\eta Rv_t.
\]
因此
\[
 \boxed{v_t=\frac{2(\rho_s-\rho_f)gR^2}{9\eta}}.
\]
该式要求 $\rho_s>\rho_f$ 且雷诺数足够小；若小球上浮，速度方向反向，大小用密度差绝对值。
解出后可以代回 $\mathrm{Re}=2\rho_fR|v_t|/\eta$ 检查是否远小于1。
这里特征长度取球直径 $2R$；若检验不满足，就必须更换阻力模型，而不是继续相信代数上求出的速度。
\end{xesolution}

\section{表面张力}\label{ch13:sec:surface}

表面张力 $\gamma$ 既可理解为沿界面单位长度的拉力，也可理解为增加单位面积界面所需的可逆功。
曲面两侧压强差满足拉普拉斯公式；球形液滴只有一个界面，肥皂泡有内外两个界面：
\begin{equation}
 \Delta p_{\rm drop}=\frac{2\gamma}{R},\qquad
 \Delta p_{\rm bubble}=\frac{4\gamma}{R}.
\end{equation}
半径越小，维持曲率所需压强差越大。毛细管中接触角为 $\theta$ 时，液面升降高度为
\begin{equation}
 h=\frac{2\gamma\cos\theta}{\rho gr}.
\end{equation}

\subsection{同一个系数为什么既是力除以长度，又是功除以面积}
设薄液膜有一条可滑动的直边，长度为 $\ell$，准静态向外移动 $\dd x$。液膜有前后两个界面，面积总增量为 $\dd A=2\ell\dd x$，外力做功为 $F\dd x$。在温度、组成不变，且过程可逆的条件下，$F\dd x=\gamma\dd A$，所以 $F=2\gamma\ell$。若只有一个界面，相应的因子2消失。
因此 $\mathrm{N/m}=\mathrm{J/m^2}$ 并非偶然的单位巧合，而是力做功与面积变化之间的关系。液膜沿界面拉紧边界，不是处处沿固定空间方向施力。

\subsection{球面压差与毛细升降的推导}
先看单界面的球形液滴。令半径准静态增加 $\dd R$，体积和面积的变化为
\[
 \dd V=4\pi R^2\dd R,\qquad \dd A=8\pi R\dd R.
\]
内外压强差做功为 $(p_{\rm in}-p_{\rm out})\dd V$，用于增加界面自由能 $\gamma\dd A$。约去公共因子后，$p_{\rm in}-p_{\rm out}=2\gamma/R$。薄壁肥皂泡有两个近似同半径的界面，面积变化加倍，压差才成为 $4\gamma/R$；“两个界面”不是把半径也乘2。

再看半径为 $r$ 的竖直毛细管，管外液面宽广。接触角 $\theta$ 从液体内部量取，圆形接触线长 $2\pi r$，界面张力的竖直合力为 $2\pi r\gamma\cos\theta$。在弯月面体积相对液柱体积可忽略、气体密度可忽略的近似下，它平衡高出外液面的液柱重量：
\[
 2\pi r\gamma\cos\theta=\rho g\pi r^2h.
\]
于是得到前式。$h$ 是相对管外液面的有符号高度：润湿时 $\theta<\pi/2$、$h>0$；不润湿时 $\theta>\pi/2$、$h<0$，表示液面下降。忽略弯月面体积通常要求管半径远小于毛细长度 $\sqrt{\gamma/(\rho g)}$。接触角变化、管壁污染或液柱尚在运动时，不能原样套用静态结果。

\Needspace{16\baselineskip}\begin{origexamplebox}{液滴合并时的界面能变化}\label{ch13:ex:teaching-droplets}
两滴同种不可压缩液体均为半径 $R$ 的球形液滴，合并后形成一个球形液滴。设表面张力为常量 $\gamma$，忽略重力。求合并后的半径、界面能减少量及新液滴的内外压差。界面能的减少是否等于总能量的消失？

\ThQDroplets
\end{origexamplebox}
\Needspace{6\baselineskip}\begin{xesolution}
\ThADroplets

先守恒的是体积而不是表面积。设新半径为 $R_f$，则
\[
 \frac43\pi R_f^3=2\frac43\pi R^3,
 \qquad R_f=2^{1/3}R.
\]
初、末总表面积分别为 $8\pi R^2$ 和 $4\pi 2^{2/3}R^2$。因此界面能减少
\[
 \Delta E_s=4\pi\gamma R^2(2-2^{2/3})>0.
\]
这部分能量可先转化为液滴内部流动、形变振动的能量，再由黏性耗散转化为内能；没有凭空消失。若过程中温度明显改变，常量 $\gamma$ 的近似还须重查。
新液滴压差为 $\Delta p_f=2\gamma/(2^{1/3}R)$，比原压差小。半径增大、曲率减小与压差下降一致；$\gamma R^2$ 的量纲是能量，$\gamma/R$ 的量纲是压强。
\end{xesolution}

\section{流体运动方程（选读）}\label{ch13:sec:equations-reading}
\begingroup\tcbset{breakable/.default=true}\begin{starread}{欧拉方程与纳维--斯托克斯方程}
无黏流体的局部动力学方程为
\[
 \rho\left(\frac{\partial\vect v}{\partial t}+\vect v\cdot\nabla\vect v\right)
 =-\nabla p+\rho\vect g.
\]
加入牛顿黏性应力后得到不可压缩纳维--斯托克斯方程（Navier--Stokes equation）
\[
 \rho\left(\frac{\partial\vect v}{\partial t}+\vect v\cdot\nabla\vect v\right)
 =-\nabla p+\eta\nabla^2\vect v+\rho\vect g,
 \qquad \nabla\cdot\vect v=0.
\]
非线性对流项使涡旋、边界层和湍流问题远比一维管流复杂。竞赛计算若采用伯努利或泊肃叶模型，
必须先说明忽略或保留了哪些项。

对可压缩牛顿流体，黏性应力的常用形式为
\[
 \boldsymbol\tau=\eta\left[\nabla\vect v+(\nabla\vect v)^{\mathsf T}
 -\frac23(\nabla\cdot\vect v)\mathbf 1\right]+\zeta(\nabla\cdot\vect v)\mathbf 1,
\]
其中 $\eta$ 是剪切黏度，$\zeta$ 是体黏度。不可压缩且 $\eta$ 均匀时才简化为动量方程中的 $\eta\nabla^2\vect v$。这说明泊肃叶流、低雷诺数绕流和可压缩声流采用的黏性模型虽同源，保留的项并不相同。
其来由是把总应力分为 $\boldsymbol\sigma=-p\mathbf I+\boldsymbol\tau$。
不含局部力偶的微小流体立方体满足角动量平衡，要求 $\tau_{xy}=\tau_{yx}$ 等，因此应力张量对称。
形变率还可分成无迹部分与纯体积变化：
\[
 \mathbf D=\left(\mathbf D-\frac13\nabla\cdot\vect v\,\mathbf I\right)
 +\frac13\nabla\cdot\vect v\,\mathbf I.
\]
各向同性牛顿流体的应力分别与这两部分成正比，所以出现剪切黏度和体黏度两个系数。
热力学压强 $p$ 与总正应力的平均值在可压缩黏性流动中未必相同，不能把包含体黏度项的全部 $\boldsymbol\tau$ 都称为无迹偏应力。
对任意小控制体，表面应力的合力密度为 $\nabla\cdot\boldsymbol\sigma$，故一般动量式为
\[
 \rho\frac{\mathrm D\vect v}{\mathrm Dt}=-\nabla p+\nabla\cdot\boldsymbol\tau+\rho\vect g.
\]
若 $\eta,\zeta$ 均匀，逐项求散度得到
\[
 \nabla\cdot\boldsymbol\tau=\eta\nabla^2\vect v+
 \left(\zeta+\frac\eta3\right)\nabla(\nabla\cdot\vect v).
\]
再用不可压缩条件，后一项才消失。方程虽然形式紧凑，三维非线性流动的存在性、正则性以及湍流描述仍涉及深刻数学问题；写出方程不等于已经求出了任意流动。
\end{starread}\endgroup

\newcommand{\FluidCoupledHgDiagram}{%
\begin{center}\begin{minipage}{\linewidth}\centering\begin{tikzpicture}[scale=.66,>=Latex]
 \fill[TextGray!30] (-.16,2.4)--(.16,2.4)--(.16,.16)--(.8,.16)--(.8,-.16)--(-.16,-.16)--cycle;
 \fill[TextGray!30] (10.8,-.16)--(14.56,-.16)--(14.56,2.4)--(14.24,2.4)--(14.24,.16)--(10.8,.16)--cycle;
 \draw[thick] (-.16,3.1)--(-.16,-.16)--(14.56,-.16)--(14.56,3.1);
 \draw[thick] (.16,3.1)--(.16,.16)--(14.24,.16)--(14.24,3.1);
 \draw[DeepBlue,thick] (-.16,2.4)--(.16,2.4) (14.24,2.4)--(14.56,2.4) (.8,-.16)--(.8,.16) (10.8,-.16)--(10.8,.16);
 \node at (5.8,.65) {封闭空气};
 \node[above right] at (.35,2.6) {$p_0$};\node[above left] at (14.05,2.6) {$p_0$};
 \draw[<->,TextGray] (-.9,0)--(-.9,2.4) node[midway,left]{$24\,\mathrm{cm}$};
 \draw[<->,TextGray] (15.2,0)--(15.2,2.4) node[midway,right]{$24\,\mathrm{cm}$};
 \foreach \xx in {0,.8,10.8,14.4}{\draw[dotted,TextGray] (\xx,-.22)--(\xx,-1.15);}
 \draw[<->,TextGray] (0,-.75)--(.8,-.75) node[midway,below=6pt]{$8\,\mathrm{cm}$};
 \draw[<->,TextGray] (.8,-.75)--(10.8,-.75) node[midway,below=6pt]{$\ell=100\,\mathrm{cm}$};
 \draw[<->,TextGray] (10.8,-.75)--(14.4,-.75) node[midway,below=6pt]{$36\,\mathrm{cm}$};
\end{tikzpicture}\figcap{平衡时两段水银柱之间封住空气。尺寸按管道中心线计，左右开口均通大气。}\label{ch13:fig:coupled-hg}\end{minipage}\end{center}}

\newcommand{\FluidCouetteDiagram}{%
\begin{center}\begin{minipage}{\linewidth}\centering\begin{tikzpicture}[scale=.9,>=Latex]
 \draw[thick] (0,0) circle(2.2) (0,0) circle(.85);
 \fill (0,0) circle(1.5pt) node[below left]{$O$};
 \draw[->,DeepBlue,thick] ({1.12*cos(20)},{1.12*sin(20)}) arc(20:85:1.12);
 \node[right] at (.65,1.22) {$\omega_0$};
 \draw[<->,TextGray] (0,0)--(-.85,0);
 \draw[TextGray,thin] (-.5,0)--(-1.1,-.5);
 \node[left] at (-1.1,-.5) {$R_1$};
 \draw[<->,TextGray] (0,0)--({2.2*cos(125)},{2.2*sin(125)}) node[midway,left=5pt]{$R_2$};
 \draw[<->,TextGray] (0,0)--(1.45,0);
 \node at (1.2,-.45) {$r$};
 \fill (1.45,0) circle(1.5pt) node[below right]{$P$};
 \draw[->,DeepBlue,thick] (1.45,0)--(1.45,.95) node[right]{$v_\theta$};
 \node[right] at (2.6,0) {外筒固定};
 \node[below] at (0,-2.55) {$z$ 轴竖直向上，图为水平截面};
\end{tikzpicture}\figcap{同轴圆筒的水平截面。内筒以 $\omega_0$ 转动，外筒固定；$P$ 为间隙中的考察点。}\label{ch13:fig:couette}\end{minipage}\end{center}}

\newcommand{\FluidHangarDiagram}{%
\begin{center}\begin{minipage}{\linewidth}\centering\begin{tikzpicture}[scale=.85,>=Latex]
 \draw[Blue,thick] plot[smooth] coordinates {(-4.000000,0.357197) (-3.955556,0.360639) (-3.911111,0.364264) (-3.866667,0.368088) (-3.822222,0.372125) (-3.777778,0.376394) (-3.733333,0.380912) (-3.688889,0.385701) (-3.644444,0.390785) (-3.600000,0.396189) (-3.555556,0.401943) (-3.511111,0.408078) (-3.466667,0.414633) (-3.422222,0.421646) (-3.377778,0.429165) (-3.333333,0.437240) (-3.288889,0.445931) (-3.244444,0.455303) (-3.200000,0.465432) (-3.155556,0.476401) (-3.111111,0.488308) (-3.066667,0.501261) (-3.022222,0.515383) (-2.977778,0.530813) (-2.933333,0.547705) (-2.888889,0.566233) (-2.844444,0.586584) (-2.800000,0.608962) (-2.755556,0.633578) (-2.711111,0.660646) (-2.666667,0.690366) (-2.622222,0.722907) (-2.577778,0.758385) (-2.533333,0.796837) (-2.488889,0.838195) (-2.444444,0.882276) (-2.400000,0.928775) (-2.355556,0.977284) (-2.311111,1.027322) (-2.266667,1.078380) (-2.222222,1.129955) (-2.177778,1.181586) (-2.133333,1.232872) (-2.088889,1.283486) (-2.044444,1.333169) (-2.000000,1.381725) (-1.955556,1.429015) (-1.911111,1.474943) (-1.866667,1.519449) (-1.822222,1.562501) (-1.777778,1.604088) (-1.733333,1.644211) (-1.688889,1.682886) (-1.644444,1.720133) (-1.600000,1.755979) (-1.555556,1.790455) (-1.511111,1.823591) (-1.466667,1.855421) (-1.422222,1.885976) (-1.377778,1.915288) (-1.333333,1.943390) (-1.288889,1.970311) (-1.244444,1.996081) (-1.200000,2.020726) (-1.155556,2.044275) (-1.111111,2.066750) (-1.066667,2.088177) (-1.022222,2.108577) (-0.977778,2.127972) (-0.933333,2.146380) (-0.888889,2.163821) (-0.844444,2.180311) (-0.800000,2.195867) (-0.755556,2.210503) (-0.711111,2.224234) (-0.666667,2.237072) (-0.622222,2.249029) (-0.577778,2.260117) (-0.533333,2.270344) (-0.488889,2.279721) (-0.444444,2.288256) (-0.400000,2.295955) (-0.355556,2.302827) (-0.311111,2.308876) (-0.266667,2.314109) (-0.222222,2.318528) (-0.177778,2.322140) (-0.133333,2.324945) (-0.088889,2.326947) (-0.044444,2.328148) (0.000000,2.328548) (0.044444,2.328148) (0.088889,2.326947) (0.133333,2.324945) (0.177778,2.322140) (0.222222,2.318528) (0.266667,2.314109) (0.311111,2.308876) (0.355556,2.302827) (0.400000,2.295955) (0.444444,2.288256) (0.488889,2.279721) (0.533333,2.270344) (0.577778,2.260117) (0.622222,2.249029) (0.666667,2.237072) (0.711111,2.224234) (0.755556,2.210503) (0.800000,2.195867) (0.844444,2.180311) (0.888889,2.163821) (0.933333,2.146380) (0.977778,2.127972) (1.022222,2.108577) (1.066667,2.088177) (1.111111,2.066750) (1.155556,2.044275) (1.200000,2.020726) (1.244444,1.996081) (1.288889,1.970311) (1.333333,1.943390) (1.377778,1.915288) (1.422222,1.885976) (1.466667,1.855421) (1.511111,1.823591) (1.555556,1.790455) (1.600000,1.755979) (1.644444,1.720133) (1.688889,1.682886) (1.733333,1.644211) (1.777778,1.604088) (1.822222,1.562501) (1.866667,1.519449) (1.911111,1.474943) (1.955556,1.429015) (2.000000,1.381725) (2.044444,1.333169) (2.088889,1.283486) (2.133333,1.232872) (2.177778,1.181586) (2.222222,1.129955) (2.266667,1.078380) (2.311111,1.027322) (2.355556,0.977284) (2.400000,0.928775) (2.444444,0.882276) (2.488889,0.838195) (2.533333,0.796837) (2.577778,0.758385) (2.622222,0.722907) (2.666667,0.690366) (2.711111,0.660646) (2.755556,0.633578) (2.800000,0.608962) (2.844444,0.586584) (2.888889,0.566233) (2.933333,0.547705) (2.977778,0.530813) (3.022222,0.515383) (3.066667,0.501261) (3.111111,0.488308) (3.155556,0.476401) (3.200000,0.465432) (3.244444,0.455303) (3.288889,0.445931) (3.333333,0.437240) (3.377778,0.429165) (3.422222,0.421646) (3.466667,0.414633) (3.511111,0.408078) (3.555556,0.401943) (3.600000,0.396189) (3.644444,0.390785) (3.688889,0.385701) (3.733333,0.380912) (3.777778,0.376394) (3.822222,0.372125) (3.866667,0.368088) (3.911111,0.364264) (3.955556,0.360639) (4.000000,0.357197)};
\draw[Blue,thick,->] (-2.444444,0.882276)--(-2.711111,0.660646);
\draw[Blue,thick] plot[smooth] coordinates {(-4.000000,0.912334) (-3.955556,0.919978) (-3.911111,0.927966) (-3.866667,0.936318) (-3.822222,0.945053) (-3.777778,0.954194) (-3.733333,0.963765) (-3.688889,0.973790) (-3.644444,0.984296) (-3.600000,0.995310) (-3.555556,1.006861) (-3.511111,1.018979) (-3.466667,1.031695) (-3.422222,1.045042) (-3.377778,1.059051) (-3.333333,1.073757) (-3.288889,1.089193) (-3.244444,1.105393) (-3.200000,1.122390) (-3.155556,1.140213) (-3.111111,1.158893) (-3.066667,1.178455) (-3.022222,1.198921) (-2.977778,1.220308) (-2.933333,1.242627) (-2.888889,1.265881) (-2.844444,1.290064) (-2.800000,1.315164) (-2.755556,1.341157) (-2.711111,1.368008) (-2.666667,1.395672) (-2.622222,1.424096) (-2.577778,1.453213) (-2.533333,1.482953) (-2.488889,1.513233) (-2.444444,1.543967) (-2.400000,1.575067) (-2.355556,1.606438) (-2.311111,1.637990) (-2.266667,1.669629) (-2.222222,1.701267) (-2.177778,1.732819) (-2.133333,1.764206) (-2.088889,1.795353) (-2.044444,1.826193) (-2.000000,1.856662) (-1.955556,1.886706) (-1.911111,1.916274) (-1.866667,1.945325) (-1.822222,1.973819) (-1.777778,2.001725) (-1.733333,2.029014) (-1.688889,2.055662) (-1.644444,2.081650) (-1.600000,2.106962) (-1.555556,2.131584) (-1.511111,2.155506) (-1.466667,2.178720) (-1.422222,2.201219) (-1.377778,2.223000) (-1.333333,2.244059) (-1.288889,2.264395) (-1.244444,2.284007) (-1.200000,2.302895) (-1.155556,2.321062) (-1.111111,2.338507) (-1.066667,2.355235) (-1.022222,2.371246) (-0.977778,2.386544) (-0.933333,2.401133) (-0.888889,2.415015) (-0.844444,2.428194) (-0.800000,2.440672) (-0.755556,2.452455) (-0.711111,2.463543) (-0.666667,2.473942) (-0.622222,2.483654) (-0.577778,2.492683) (-0.533333,2.501030) (-0.488889,2.508699) (-0.444444,2.515692) (-0.400000,2.522012) (-0.355556,2.527661) (-0.311111,2.532640) (-0.266667,2.536952) (-0.222222,2.540598) (-0.177778,2.543580) (-0.133333,2.545898) (-0.088889,2.547553) (-0.044444,2.548545) (0.000000,2.548876) (0.044444,2.548545) (0.088889,2.547553) (0.133333,2.545898) (0.177778,2.543580) (0.222222,2.540598) (0.266667,2.536952) (0.311111,2.532640) (0.355556,2.527661) (0.400000,2.522012) (0.444444,2.515692) (0.488889,2.508699) (0.533333,2.501030) (0.577778,2.492683) (0.622222,2.483654) (0.666667,2.473942) (0.711111,2.463543) (0.755556,2.452455) (0.800000,2.440672) (0.844444,2.428194) (0.888889,2.415015) (0.933333,2.401133) (0.977778,2.386544) (1.022222,2.371246) (1.066667,2.355235) (1.111111,2.338507) (1.155556,2.321062) (1.200000,2.302895) (1.244444,2.284007) (1.288889,2.264395) (1.333333,2.244059) (1.377778,2.223000) (1.422222,2.201219) (1.466667,2.178720) (1.511111,2.155506) (1.555556,2.131584) (1.600000,2.106962) (1.644444,2.081650) (1.688889,2.055662) (1.733333,2.029014) (1.777778,2.001725) (1.822222,1.973819) (1.866667,1.945325) (1.911111,1.916274) (1.955556,1.886706) (2.000000,1.856662) (2.044444,1.826193) (2.088889,1.795353) (2.133333,1.764206) (2.177778,1.732819) (2.222222,1.701267) (2.266667,1.669629) (2.311111,1.637990) (2.355556,1.606438) (2.400000,1.575067) (2.444444,1.543967) (2.488889,1.513233) (2.533333,1.482953) (2.577778,1.453213) (2.622222,1.424096) (2.666667,1.395672) (2.711111,1.368008) (2.755556,1.341157) (2.800000,1.315164) (2.844444,1.290064) (2.888889,1.265881) (2.933333,1.242627) (2.977778,1.220308) (3.022222,1.198921) (3.066667,1.178455) (3.111111,1.158893) (3.155556,1.140213) (3.200000,1.122390) (3.244444,1.105393) (3.288889,1.089193) (3.333333,1.073757) (3.377778,1.059051) (3.422222,1.045042) (3.466667,1.031695) (3.511111,1.018979) (3.555556,1.006861) (3.600000,0.995310) (3.644444,0.984296) (3.688889,0.973790) (3.733333,0.963765) (3.777778,0.954194) (3.822222,0.945053) (3.866667,0.936318) (3.911111,0.927966) (3.955556,0.919978) (4.000000,0.912334)};
\draw[Blue,thick,->] (-2.444444,1.543967)--(-2.711111,1.368008);
\draw[Blue,thick] plot[smooth] coordinates {(-4.000000,1.497260) (-3.955556,1.507103) (-3.911111,1.517277) (-3.866667,1.527793) (-3.822222,1.538663) (-3.777778,1.549898) (-3.733333,1.561508) (-3.688889,1.573505) (-3.644444,1.585897) (-3.600000,1.598697) (-3.555556,1.611912) (-3.511111,1.625551) (-3.466667,1.639623) (-3.422222,1.654135) (-3.377778,1.669092) (-3.333333,1.684499) (-3.288889,1.700360) (-3.244444,1.716675) (-3.200000,1.733445) (-3.155556,1.750667) (-3.111111,1.768338) (-3.066667,1.786450) (-3.022222,1.804994) (-2.977778,1.823960) (-2.933333,1.843334) (-2.888889,1.863099) (-2.844444,1.883236) (-2.800000,1.903724) (-2.755556,1.924539) (-2.711111,1.945656) (-2.666667,1.967047) (-2.622222,1.988681) (-2.577778,2.010528) (-2.533333,2.032554) (-2.488889,2.054726) (-2.444444,2.077010) (-2.400000,2.099370) (-2.355556,2.121772) (-2.311111,2.144179) (-2.266667,2.166559) (-2.222222,2.188876) (-2.177778,2.211096) (-2.133333,2.233189) (-2.088889,2.255123) (-2.044444,2.276867) (-2.000000,2.298394) (-1.955556,2.319676) (-1.911111,2.340688) (-1.866667,2.361407) (-1.822222,2.381809) (-1.777778,2.401874) (-1.733333,2.421582) (-1.688889,2.440916) (-1.644444,2.459860) (-1.600000,2.478398) (-1.555556,2.496516) (-1.511111,2.514202) (-1.466667,2.531444) (-1.422222,2.548233) (-1.377778,2.564558) (-1.333333,2.580411) (-1.288889,2.595785) (-1.244444,2.610673) (-1.200000,2.625069) (-1.155556,2.638968) (-1.111111,2.652365) (-1.066667,2.665255) (-1.022222,2.677636) (-0.977778,2.689503) (-0.933333,2.700854) (-0.888889,2.711687) (-0.844444,2.722000) (-0.800000,2.731790) (-0.755556,2.741056) (-0.711111,2.749797) (-0.666667,2.758012) (-0.622222,2.765700) (-0.577778,2.772859) (-0.533333,2.779490) (-0.488889,2.785592) (-0.444444,2.791164) (-0.400000,2.796206) (-0.355556,2.800718) (-0.311111,2.804700) (-0.266667,2.808151) (-0.222222,2.811071) (-0.177778,2.813460) (-0.133333,2.815319) (-0.088889,2.816646) (-0.044444,2.817443) (0.000000,2.817708) (0.044444,2.817443) (0.088889,2.816646) (0.133333,2.815319) (0.177778,2.813460) (0.222222,2.811071) (0.266667,2.808151) (0.311111,2.804700) (0.355556,2.800718) (0.400000,2.796206) (0.444444,2.791164) (0.488889,2.785592) (0.533333,2.779490) (0.577778,2.772859) (0.622222,2.765700) (0.666667,2.758012) (0.711111,2.749797) (0.755556,2.741056) (0.800000,2.731790) (0.844444,2.722000) (0.888889,2.711687) (0.933333,2.700854) (0.977778,2.689503) (1.022222,2.677636) (1.066667,2.665255) (1.111111,2.652365) (1.155556,2.638968) (1.200000,2.625069) (1.244444,2.610673) (1.288889,2.595785) (1.333333,2.580411) (1.377778,2.564558) (1.422222,2.548233) (1.466667,2.531444) (1.511111,2.514202) (1.555556,2.496516) (1.600000,2.478398) (1.644444,2.459860) (1.688889,2.440916) (1.733333,2.421582) (1.777778,2.401874) (1.822222,2.381809) (1.866667,2.361407) (1.911111,2.340688) (1.955556,2.319676) (2.000000,2.298394) (2.044444,2.276867) (2.088889,2.255123) (2.133333,2.233189) (2.177778,2.211096) (2.222222,2.188876) (2.266667,2.166559) (2.311111,2.144179) (2.355556,2.121772) (2.400000,2.099370) (2.444444,2.077010) (2.488889,2.054726) (2.533333,2.032554) (2.577778,2.010528) (2.622222,1.988681) (2.666667,1.967047) (2.711111,1.945656) (2.755556,1.924539) (2.800000,1.903724) (2.844444,1.883236) (2.888889,1.863099) (2.933333,1.843334) (2.977778,1.823960) (3.022222,1.804994) (3.066667,1.786450) (3.111111,1.768338) (3.155556,1.750667) (3.200000,1.733445) (3.244444,1.716675) (3.288889,1.700360) (3.333333,1.684499) (3.377778,1.669092) (3.422222,1.654135) (3.466667,1.639623) (3.511111,1.625551) (3.555556,1.611912) (3.600000,1.598697) (3.644444,1.585897) (3.688889,1.573505) (3.733333,1.561508) (3.777778,1.549898) (3.822222,1.538663) (3.866667,1.527793) (3.911111,1.517277) (3.955556,1.507103) (4.000000,1.497260)};
\draw[Blue,thick,->] (-2.444444,2.077010)--(-2.711111,1.945656);
\draw[Blue,thick] plot[smooth] coordinates {(-4.000000,2.155276) (-3.955556,2.165382) (-3.911111,2.175721) (-3.866667,2.186297) (-3.822222,2.197111) (-3.777778,2.208165) (-3.733333,2.219461) (-3.688889,2.231000) (-3.644444,2.242783) (-3.600000,2.254809) (-3.555556,2.267079) (-3.511111,2.279591) (-3.466667,2.292343) (-3.422222,2.305333) (-3.377778,2.318559) (-3.333333,2.332016) (-3.288889,2.345699) (-3.244444,2.359605) (-3.200000,2.373725) (-3.155556,2.388055) (-3.111111,2.402585) (-3.066667,2.417308) (-3.022222,2.432213) (-2.977778,2.447292) (-2.933333,2.462534) (-2.888889,2.477926) (-2.844444,2.493456) (-2.800000,2.509112) (-2.755556,2.524880) (-2.711111,2.540746) (-2.666667,2.556696) (-2.622222,2.572714) (-2.577778,2.588785) (-2.533333,2.604893) (-2.488889,2.621022) (-2.444444,2.637156) (-2.400000,2.653280) (-2.355556,2.669376) (-2.311111,2.685428) (-2.266667,2.701420) (-2.222222,2.717335) (-2.177778,2.733159) (-2.133333,2.748876) (-2.088889,2.764469) (-2.044444,2.779924) (-2.000000,2.795227) (-1.955556,2.810362) (-1.911111,2.825317) (-1.866667,2.840078) (-1.822222,2.854631) (-1.777778,2.868966) (-1.733333,2.883068) (-1.688889,2.896928) (-1.644444,2.910534) (-1.600000,2.923876) (-1.555556,2.936944) (-1.511111,2.949729) (-1.466667,2.962221) (-1.422222,2.974412) (-1.377778,2.986295) (-1.333333,2.997861) (-1.288889,3.009103) (-1.244444,3.020015) (-1.200000,3.030591) (-1.155556,3.040824) (-1.111111,3.050709) (-1.066667,3.060242) (-1.022222,3.069416) (-0.977778,3.078229) (-0.933333,3.086674) (-0.888889,3.094750) (-0.844444,3.102451) (-0.800000,3.109776) (-0.755556,3.116719) (-0.711111,3.123280) (-0.666667,3.129455) (-0.622222,3.135242) (-0.577778,3.140638) (-0.533333,3.145643) (-0.488889,3.150253) (-0.444444,3.154467) (-0.400000,3.158284) (-0.355556,3.161703) (-0.311111,3.164722) (-0.266667,3.167340) (-0.222222,3.169557) (-0.177778,3.171372) (-0.133333,3.172785) (-0.088889,3.173794) (-0.044444,3.174399) (0.000000,3.174601) (0.044444,3.174399) (0.088889,3.173794) (0.133333,3.172785) (0.177778,3.171372) (0.222222,3.169557) (0.266667,3.167340) (0.311111,3.164722) (0.355556,3.161703) (0.400000,3.158284) (0.444444,3.154467) (0.488889,3.150253) (0.533333,3.145643) (0.577778,3.140638) (0.622222,3.135242) (0.666667,3.129455) (0.711111,3.123280) (0.755556,3.116719) (0.800000,3.109776) (0.844444,3.102451) (0.888889,3.094750) (0.933333,3.086674) (0.977778,3.078229) (1.022222,3.069416) (1.066667,3.060242) (1.111111,3.050709) (1.155556,3.040824) (1.200000,3.030591) (1.244444,3.020015) (1.288889,3.009103) (1.333333,2.997861) (1.377778,2.986295) (1.422222,2.974412) (1.466667,2.962221) (1.511111,2.949729) (1.555556,2.936944) (1.600000,2.923876) (1.644444,2.910534) (1.688889,2.896928) (1.733333,2.883068) (1.777778,2.868966) (1.822222,2.854631) (1.866667,2.840078) (1.911111,2.825317) (1.955556,2.810362) (2.000000,2.795227) (2.044444,2.779924) (2.088889,2.764469) (2.133333,2.748876) (2.177778,2.733159) (2.222222,2.717335) (2.266667,2.701420) (2.311111,2.685428) (2.355556,2.669376) (2.400000,2.653280) (2.444444,2.637156) (2.488889,2.621022) (2.533333,2.604893) (2.577778,2.588785) (2.622222,2.572714) (2.666667,2.556696) (2.711111,2.540746) (2.755556,2.524880) (2.800000,2.509112) (2.844444,2.493456) (2.888889,2.477926) (2.933333,2.462534) (2.977778,2.447292) (3.022222,2.432213) (3.066667,2.417308) (3.111111,2.402585) (3.155556,2.388055) (3.200000,2.373725) (3.244444,2.359605) (3.288889,2.345699) (3.333333,2.332016) (3.377778,2.318559) (3.422222,2.305333) (3.466667,2.292343) (3.511111,2.279591) (3.555556,2.267079) (3.600000,2.254809) (3.644444,2.242783) (3.688889,2.231000) (3.733333,2.219461) (3.777778,2.208165) (3.822222,2.197111) (3.866667,2.186297) (3.911111,2.175721) (3.955556,2.165382) (4.000000,2.155276)};
\draw[Blue,thick,->] (-2.444444,2.637156)--(-2.711111,2.540746);

 \draw[thick] (-2.2,0) arc(180:0:2.2);
 \draw[->,TextGray] (-4.2,0)--(4.5,0) node[right]{$x$};
 \fill (-2.2,0) circle(1.6pt);\node[below] at (-2.2,-.15) {$A$（门口）};
 \fill (0,0) circle(1.6pt) node[below]{$O$};
 \coordinate (B) at ({2.2*cos(55)},{2.2*sin(55)});
 \draw[TextGray] (0,0)--(B);\fill (B) circle(1.6pt);
 \draw[TextGray,thin] (.8,1.7)--(B);
 \node[left] at (.7,1.7) {$B$};\node[left] at (.55,.95) {$R$};
 \draw[TextGray] (.65,0) arc(0:55:.65);\node at (.92,.35) {$\theta$};
 \node[above left] at (-2.4,4.1) {来流向左};
\end{tikzpicture}\figcap{半圆柱机库的横截面与题设流线。$A$ 为敞开的门口，$R$ 为半径，$\theta$ 从正 $x$ 方向量起。}\label{ch13:fig:hangar}\end{minipage}\end{center}}

\newcommand{\FluidFunnelDiagram}{%
\begin{center}\begin{minipage}{\linewidth}\centering\begin{tikzpicture}[scale=.85,>=Latex]
 \def\aa{.16}\def\RR{2.4}\def\HH{3.4}\def\hh{2.65}
 \pgfmathsetmacro{\rh}{\aa+(\RR-\aa)*\hh/\HH}
 \fill[Blue!12] (-\aa,0)--(-\rh,\hh)--(\rh,\hh)--(\aa,0)--cycle;
 \draw[thick] (-\RR,\HH)--(-\aa,0) (\aa,0)--(\RR,\HH);
 \draw[DeepBlue,thick] (-\rh,\hh)--(\rh,\hh);
 \draw[DeepBlue,thick] (-\aa,0)--(-\aa,-1) (\aa,0)--(\aa,-1);
 \draw[->,DeepBlue,thick] (0,-.4)--(0,-1.25) node[midway,right=7pt]{$Q$};
 \draw[<->,TextGray] (-3,0)--(-3,\HH) node[midway,left]{$H$};
 \draw[<->,TextGray] (2.9,0)--(2.9,\hh) node[midway,right]{$h$};
 \draw[dotted,TextGray] (\aa,0)--(2.9,0) (\rh,\hh)--(2.9,\hh);
 \draw[<->,TextGray] (0,3.85)--(\RR,3.85) node[midway,above]{$R$};
 \draw[TextGray,thin] (-\aa,-1.25)--(-\aa,-1.68) (\aa,-1.25)--(\aa,-1.68);
 \draw[->,TextGray] (-.9,-1.55)--(-\aa,-1.55);
 \draw[->,TextGray] (.9,-1.55)--(\aa,-1.55);
 \node[below] at (0,-1.65) {有限开口 $2a$};
\end{tikzpicture}\figcap{圆锥形漏斗与连通的有限出口。$H$、$h$ 均从出口量起，图示排液过程中的一个状态。}\label{ch13:fig:funnel}\end{minipage}\end{center}}

\newcommand{\FluidShoreDiagram}{%
\begin{center}\begin{minipage}{\linewidth}\centering\begin{tikzpicture}[scale=.9,>=Latex]
 \begin{scope}[shift={(-5.6,1.4)}]
  \fill[Blue!12] (0,0)--(4.2,0)--(4.2,-1.2)--cycle;
  \draw[DeepBlue,thick] (0,0)--(4.2,0);
  \draw[Ink,thick] (-.7,.2)--(4.2,-1.2);
  \draw[TextGray] (1.1,0) arc(0:-15.9454:1.1);
  \node at (1.4,-.22) {$\alpha$};
  \fill (0,0) circle(1.5pt) node[above] {岸};
  \node[above] at (3.1,.15) {水面};
  \node[below] at (3.1,-1.3) {岸坡剖面};
 \end{scope}
 \begin{scope}[shift={(.2,0)}]
  \fill[Blue!8] (0,0) rectangle (3.6,2.7);
  \draw[->,TextGray] (0,0)--(3.9,0) node[right]{$x$};
  \draw[->,TextGray] (0,0)--(0,3) node[above]{$y$};
  \draw[DeepBlue,very thick,domain=0:180,samples=120,variable=\p]
   plot ({pi*\p/180+sin(\p)},{1+cos(\p)});
  \fill (0,2) circle(1.8pt) node[left]{$B(0,y_0)$};
  \fill (pi,0) circle(1.8pt) node[below]{$A$};
  \draw[->,DeepBlue,thick] (0,2)--(.8,2) node[midway,above=5pt]{$\vect c_0$};
  \node[below] at (1,-.4) {岸线 $y=0$};
 \end{scope}
\end{tikzpicture}\figcap{左为岸坡剖面，水只在水面与坡底之间；右为原题所给的弯向岸边的射线示意，$B$ 在离岸坐标轴上。}\label{ch13:fig:shore}\end{minipage}\end{center}}

\newcommand{\FluidHoseDiagram}{%
\begin{center}\begin{minipage}{\linewidth}\centering\begin{tikzpicture}[scale=.95,>=Latex]
 \def\ang{25}\def\LL{2.65}
 \coordinate (O) at (0,0);
 \coordinate (E) at ({-\LL*sin(\ang)},{-\LL*cos(\ang)});
 \coordinate (P) at ({\LL*(cos(\ang)-sin(\ang))},{-\LL*(cos(\ang)+sin(\ang))});
 \draw[dashed,TextGray] (O)--(0,-\LL)--(\LL,-\LL);
 \draw[DeepBlue,line width=2pt] (O)--(E)--(P);
 \foreach \pt in {O,E,P}{\fill (\pt) circle(1.7pt);}
 \node[above] at (O) {$O$};\node[left=5pt] at (E) {$E$};\node[below=7pt] at (P) {$P$};
 \draw[->,Ink] ($(O)!.25!(E)$)--($(O)!.65!(E)$);
 \draw[->,Ink] ($(E)!.25!(P)$)--($(E)!.65!(P)$);
 \draw[TextGray] (0,-.75) arc(-90:{-90-\ang}:.75);
 \node[right] at (.15,-.95) {$\theta$};
 \coordinate (d1) at ({-.6*cos(\ang)},{.6*sin(\ang)});
 \coordinate (d2) at ({-.55*sin(\ang)},{-.55*cos(\ang)});
 \draw[<->,TextGray] ($(O)+(d1)$)--($(E)+(d1)$) node[midway,left=5pt]{$L$};
 \draw[<->,TextGray] ($(E)+(d2)$)--($(P)+(d2)$) node[midway,below=5pt]{$L$};
 \draw[->,Ink] (2.6,-.2)--(2.6,-1.4) node[right]{$\vect g$};
 \draw[->,DeepBlue,thick] (P)--++({.85*cos(\ang)},{-.85*sin(\ang)}) node[right]{出流方向};
\end{tikzpicture}\figcap{两段等长管在 $E$ 处保持直角；实线仅示意未知偏角，虚线为题设的角度参考构形。}\label{ch13:fig:hose}\end{minipage}\end{center}}

\newcommand{\FluidHoseAnalysisDiagram}{%
\begin{center}\begin{minipage}{\linewidth}\centering\begin{tikzpicture}[scale=.95,>=Latex]
 \def\ang{31.0809622}\def\LL{2.65}
 \coordinate (O) at (0,0);
 \coordinate (E) at ({-\LL*sin(\ang)},{-\LL*cos(\ang)});
 \coordinate (P) at ({\LL*(cos(\ang)-sin(\ang))},{-\LL*(cos(\ang)+sin(\ang))});
 \draw[dashed,TextGray] (O)--(0,-\LL)--(\LL,-\LL);
 \draw[DeepBlue,line width=2pt] (O)--(E)--(P);
 \foreach \pt in {O,E,P}{\fill (\pt) circle(1.7pt);}
 \node[above] at (O) {$O$};\node[left=5pt] at (E) {$E$};\node[below=7pt] at (P) {$P$};
 \draw[->,Ink] ($(O)!.25!(E)$)--($(O)!.65!(E)$);
 \draw[->,Ink] ($(E)!.25!(P)$)--($(E)!.65!(P)$);
 \draw[TextGray] (0,-.75) arc(-90:{-90-\ang}:.75);
 \node[right] at (.15,-.95) {$\theta=31.1^\circ$};
 \coordinate (d1) at ({-.6*cos(\ang)},{.6*sin(\ang)});
 \coordinate (d2) at ({-.55*sin(\ang)},{-.55*cos(\ang)});
 \draw[<->,TextGray] ($(O)+(d1)$)--($(E)+(d1)$) node[midway,left=5pt]{$L$};
 \draw[<->,TextGray] ($(E)+(d2)$)--($(P)+(d2)$) node[midway,below=5pt]{$L$};
 \draw[->,Ink] (2.6,-.2)--(2.6,-1.4) node[right]{$\vect g$};
 \draw[->,DeepBlue,thick] (P)--++({.85*cos(\ang)},{-.85*sin(\ang)}) node[right]{$u=Q/A$};
\end{tikzpicture}\figcap{两段等长管在 $E$ 处保持直角；实线按求得的稳定角度绘制，虚线为比较用初始构形。}\label{ch13:fig:hose-analysis}\end{minipage}\end{center}}

\section{流体动力学综合问题}\label{ch13:sec:problems}

\Needspace{16\baselineskip}\begin{origexamplebox}{两段水银柱与封闭空气的耦合振动}\label{ch13:ex:coupled-hg}
等截面的 U 形管内封有两段水银柱，中间夹着一段空气柱。平衡时，两侧开口处的水银自由面均高出
水平空气柱 $h=24\,\mathrm{cm}$；左、右水银柱沿管道的总长度分别为
$L_1=24\,\mathrm{cm}+8\,\mathrm{cm}$、$L_2=24\,\mathrm{cm}+36\,\mathrm{cm}$，
空气柱长 $\ell=100\,\mathrm{cm}$。大气压等于 $76\,\mathrm{cmHg}$，取
$g=9.8\,\mathrm{m/s^2}$。忽略阻力，把空气视为双原子理想气体并认为微振动过程绝热，求两个本征频率。

\FluidCoupledHgDiagram
\end{origexamplebox}
\Needspace{6\baselineskip}\begin{xesolution}
\ThACoupledMercury

图~\ref{ch13:fig:coupled-hg}中，水平气柱中心线到两个自由面的竖直高度均为 $h$。
平衡时空气压强等于大气压再加一段高度为 $h$ 的水银柱压强：
\[
 \bar p=p_0+\rho gh=\rho g(0.76+0.24)\,\mathrm m=\rho g\ell .
\]
这里 $\rho$ 是水银密度。定义 $q_1$、$q_2$ 分别为左、右水银内侧界面向空气柱内的位移，
即左界面向右为正，右界面向左为正；对应外侧自由面均下降同样距离。
因此空气体积的一阶相对变化为
\[
 \frac{\delta V}{V_0}=-\frac{q_1+q_2}{\ell}.
\]
绝热关系 $pV^\gamma=\text{常量}$ 给出
\[
 \delta p=\gamma\bar p\frac{q_1+q_2}{\ell}.
\]

以左水银柱为例。自由面下降 $q_1$ 使重力压头减少 $\rho gq_1$；空气压强增加 $\delta p$，也把水银柱推离空气。
沿 $q_1$ 正向写牛顿第二定律，管截面积约去后得到
\[
 L_1\ddot q_1+gq_1+\frac{\delta p}{\rho}=0.
\]
右侧同理。利用 $\bar p/(\rho\ell)=g$，两式成为
\begin{align*}
 L_1\ddot q_1+gq_1+\gamma g(q_1+q_2)&=0,\\
 L_2\ddot q_2+gq_2+\gamma g(q_1+q_2)&=0.
\end{align*}
令 $q_i=Q_i e^{\ii\omega t}$，非零振幅要求
\[
 \begin{vmatrix}
 g(1+\gamma)-\omega^2L_1&\gamma g\\
 \gamma g&g(1+\gamma)-\omega^2L_2
 \end{vmatrix}=0.
\]
记 $x=\omega^2/g$，便有
\[
 L_1L_2x^2-(1+\gamma)(L_1+L_2)x+(1+2\gamma)=0.
\]
代入 $L_1=0.32\,\mathrm m$、$L_2=0.60\,\mathrm m$、$\gamma=7/5$，得到
\[
 x_-=2.1071\,\mathrm{m^{-1}},\qquad x_+=9.3929\,\mathrm{m^{-1}},
\]
所以
\[
 \boxed{\omega_-=4.54\,\mathrm{rad/s},\quad \omega_+=9.59\,\mathrm{rad/s}},
 \qquad
 \boxed{f_-=0.723\,\mathrm{Hz},\quad f_+=1.53\,\mathrm{Hz}}.
\]
低频模中两段水银大体反向运动，空气压缩较小；高频模中两者大体同时压缩空气。
若振动足够慢、空气近似等温，只需把 $\gamma$ 改为 $1$，会得到
$f_-=0.719\,\mathrm{Hz}$、$f_+=1.36\,\mathrm{Hz}$。
\end{xesolution}

\Needspace{16\baselineskip}\begin{origexamplebox}{同轴圆筒间的黏性旋转流}\label{ch13:ex:couette}
半径为 $R_1$ 与 $R_2$（$R_2>R_1$）的两个长圆筒竖直同轴放置，间隙充满密度为 $\rho$、动力黏度为 $\eta$ 的
不可压缩牛顿流体。内筒以恒定角速度 $\omega_0$ 转动，外筒静止，流动稳定且轴对称。

（1）求角速度分布 $\omega(r)$；（2）先求同一高度的压强分布，再补出竖直方向的变化；（3）说明柱坐标中的剪切应力公式
\[
 \tau_{r\theta}=\eta\left(\frac1r\frac{\partial v_r}{\partial\theta}
 +\frac{\partial v_\theta}{\partial r}-\frac{v_\theta}{r}\right)
\]
为何与直角坐标中的形式不同。

\FluidCouetteDiagram
\end{origexamplebox}
\Needspace{6\baselineskip}\begin{xesolution}
\ThACouette

稳定轴对称流满足 $v_r=v_z=0$、$v_\theta=r\omega(r)$。
图~\ref{ch13:fig:couette}中 $P$ 的切向速度为 $v_\theta$，半径为 $r$；所画水平截面上各点的高度相同。将速度代入切应力公式，得到
\[
 \tau_{r\theta}=\eta\left(\frac{\dd v_\theta}{\dd r}-\frac{v_\theta}{r}\right)
 =\eta r\frac{\dd\omega}{\dd r}.
\]
任取半径为 $r$、长度为 $l$ 的圆柱面，黏性力矩为
$2\pi rl\,\tau_{r\theta}\,r$。稳定时各圆柱面传递的力矩相同，故
$r^2\tau_{r\theta}$ 为常量，即 $r^3\omega'(r)$ 为常量。积分并使用
$\omega(R_1)=\omega_0$、$\omega(R_2)=0$，得
\[
 \boxed{\omega(r)=\frac{\omega_0R_1^2}{R_2^2-R_1^2}
 \left(\frac{R_2^2}{r^2}-1\right)}.
\]
这一步可写为 $r^3\omega'=K$，因此 $\omega=C_1+C_2/r^2$。外壁条件给 $C_1=-C_2/R_2^2$，内壁条件再给
$C_2=\omega_0R_1^2R_2^2/(R_2^2-R_1^2)$。两边界均满足，且 $\omega_0=0$ 时全流体静止。

径向速度虽为零，流体微元仍有向心加速度 $-v_\theta^2/r$。径向欧拉方程因此给出
\[
 \frac{\dd p}{\dd r}=\rho\frac{v_\theta^2}{r}.
\]
令 $C=\omega_0R_1^2/(R_2^2-R_1^2)$，由
$v_\theta=C(R_2^2/r-r)$ 积分可得
\[
 \boxed{p(r)-p(R_1)=\rho C^2\left[
 \frac{r^2-R_1^2}{2}-2R_2^2\ln\frac r{R_1}
 -\frac{R_2^4}{2}\left(\frac1{r^2}-\frac1{R_1^2}\right)
 \right]}.
\]
因 $\dd p/\dd r=\rho v_\theta^2/r\ge0$，压强必向外增大，这也直接检查了积分式的符号。
上式只比较同一高度的两点。取竖直向上为 $z$，由于 $v_z=0$，竖直平衡还要求
$\partial p/\partial z=-\rho g$。若把方括号中的径向差记为 $\Delta p_r(r)$，完整压强场为
\[
 p(r,z)=p(R_1,z_0)+\Delta p_r(r)-\rho g(z-z_0).
\]
在不转动极限，径向差消失，只剩熟悉的静水压分布。

柱坐标基矢 $\vect e_r,\vect e_\theta$ 随方位角改变，且方位方向的实际弧长是 $r\,\dd\theta$，所以：
（i）$r^{-1}\partial v_r/\partial\theta$ 来自弧长尺度；
（ii）$-v_\theta/r$ 来自基矢方向随位置转动；
（iii）圆柱面传递力矩时还多一个力臂 $r$，因此守恒的是 $r^2\tau_{r\theta}$，而非平面流中的恒定剪切应力。
同样，曲线运动还产生直角平行流中没有的向心压强梯度。
\end{xesolution}

\Needspace{16\baselineskip}\begin{origexamplebox}{半圆柱形机库受到的风力}\label{ch13:ex:hangar}
长度为 $l$、横截面半径为 $R$ 的半圆柱形机库，处在垂直于柱轴的均匀风中。空气密度为 $\rho$，忽略重力与黏性，外部速度势为
\[
 \phi=-v_0\left(r+\frac{R^2}{r}\right)\cos\theta .
\]
地面处的门口 $A$ 敞开，使库内空气可视为静止且压强等于 $A$ 点外部压强。求风对机库曲面部分的合力。

\FluidHangarDiagram
\end{origexamplebox}
\Needspace{6\baselineskip}\begin{xesolution}
\ThAHangar

图~\ref{ch13:fig:hangar}中，$\theta$ 从正 $x$ 方向量起，门口 $A$ 对应 $\pi$。
由 $\vect v=\nabla\phi$，在曲面 $r=R$ 上有
\[
 v_r=0,\qquad v_\theta=\frac1r\frac{\partial\phi}{\partial\theta}
 =2v_0\sin\theta.
\]
外部沿流线使用伯努利方程：
\[
 p_{\rm out}(\theta)=p_\infty+\frac12\rho v_0^2
 \left(1-4\sin^2\theta\right).
\]
门口 $A$ 对应 $\theta=\pi$，是停滞点，故库内静止空气压强为
\[
 p_{\rm in}=p_A=p_\infty+\frac12\rho v_0^2.
\]
于是曲面内外压强差为
\[
 \Delta p=p_{\rm in}-p_{\rm out}=2\rho v_0^2\sin^2\theta.
\]
曲面面积元为 $\dd S=lR\,\dd\theta$，外法线
$\vect n=(\cos\theta,\sin\theta)$。左右对称使水平分量相消，而竖直分量为
\begin{align*}
 F_y&=2\rho v_0^2lR\int_0^\pi\sin^3\theta\,\dd\theta\\
 &=\boxed{\frac83\rho v_0^2Rl},\qquad F_x=0.
\end{align*}
合力竖直向上。若门关闭且库内压强不是停滞点压强，这个结论会随内部压强改变，不能直接套用。
\end{xesolution}

\Needspace{16\baselineskip}\begin{origexamplebox}{圆锥形漏斗排液}\label{ch13:ex:funnel}
一个尖端向下的圆锥形漏斗，高为 $H$、顶面半径为 $R$，初始装满水。底部有半径为 $a$ 的小孔，打开后自由流出。
忽略自由液面速度和黏性，求水面从高度 $H$ 降到 $H/2$ 所需的时间，并说明何时必须考虑黏性。

\FluidFunnelDiagram
\end{origexamplebox}
\Needspace{6\baselineskip}\begin{xesolution}
\ThAFunnel

图~\ref{ch13:fig:funnel}保留有限出口。先取 $a\ll R$ 的尖锥近似，忽略小孔截断锥尖带来的几何修正。
由相似三角形，水深为 $h$ 时的自由液面面积为
\[
 A(h)=\pi\left(\frac{Rh}{H}\right)^2.
\]
小孔处托里拆利速度为 $v=\sqrt{2gh}$，流量
$Q=\pi a^2\sqrt{2gh}$。体积守恒 $A(h)\dd h/\dd t=-Q$，故
\[
 \dd t=-\frac{R^2}{a^2H^2\sqrt{2g}}h^{3/2}\,\dd h.
\]
从 $H$ 积分到 $H/2$：
\[
 \boxed{t=\frac{2R^2}{5a^2}\sqrt{\frac{H}{2g}}
 \left(1-2^{-5/2}\right)}.
\]
量纲由 $\sqrt{H/g}$ 给出时间。这里忽略黏性的条件不是“小孔很小”，而是小孔足够短、雷诺数足够大，流出损失可用接近 $1$ 的流量系数描述。
实际锐边小孔若流量系数为 $C_d<1$，时间应除以 $C_d$；若出口是一根长细管，则黏性压降占主导，必须改用泊肃叶流模型。
若不忽略有限孔径，令 $H$ 为出口到顶面的真实高度，截面半径应为
$r(h)=a+(R-a)h/H$。设 $b=(R-a)/H$，则把同一个体积守恒式中的面积换成 $\pi(a+bh)^2$，得到
\[
 t=\frac1{\sqrt{2g}}\left[
 2\sqrt h+\frac{4b}{3a}h^{3/2}+\frac{2b^2}{5a^2}h^{5/2}
 \right]_{H/2}^{H}.
\]
当 $a/R\to0$ 时，最后一项占主导，恢复前面的尖锥结果。有限开口与理想模型的关系因此是可检验的近似，不是把管口画成封闭尖点。
\end{xesolution}

\Needspace{16\baselineskip}\begin{origexamplebox}{浅水波向岸边弯折的轨迹}\label{ch13:ex:shore}
平直海岸的水深随离岸距离线性增加。取 $x$ 轴沿岸、$y$ 轴垂直海岸并指向海面，岸线为 $y=0$，水深
$h(y)=y\tan\alpha$。浅水波速满足 $c=\sqrt{gh}=C\sqrt y$。一列波包从
$B=(0,y_0)$ 出发，初始传播方向平行于岸线。求波包射线，并证明它是一段摆线。

\FluidShoreDiagram
\end{origexamplebox}
\Needspace{6\baselineskip}\begin{xesolution}
\ThAShore

介质只随 $y$ 变化，沿岸方向平移对称。令 $\theta$ 为射线与指向岸边的法线方向（$-y$ 方向）的夹角，
图~\ref{ch13:fig:shore}中 $B$ 在 $y$ 轴上，故初始沿岸坐标为零。
射线折射不变量为
\[
 \frac{\sin\theta}{c(y)}=\text{常量}.
\]
在 $B$ 点射线平行于岸线，$\theta_0=\pi/2$，故
\[
 \sin\theta=\frac{c(y)}{c(y_0)}=\sqrt{\frac y{y_0}}.
\]
沿传播方向 $y$ 减小、$x$ 增加，因而
\[
 \frac{\dd x}{\dd y}=-\tan\theta.
\]
作参数
\[
 \varphi=\pi-2\theta,\qquad 0\le\varphi\le\pi,
\]
由 $y/y_0=\sin^2\theta=(1+\cos\varphi)/2$，有
\[
 \frac{\dd y}{\dd\varphi}=-\frac{y_0}{2}\sin\varphi,
 \qquad \tan\theta=\cot(\varphi/2).
\]
所以
\[
 \frac{\dd x}{\dd\varphi}
 =\left(-\cot\frac\varphi2\right)\left(-\frac{y_0}2\sin\varphi\right)
 =\frac{y_0}2(1+\cos\varphi).
\]
从0积分到 $\varphi$，并取 $x(0)=0$，得到
\[
 \boxed{x=\frac{y_0}{2}(\varphi+\sin\varphi)},\qquad
 \boxed{y=\frac{y_0}{2}(1+\cos\varphi)}.
\]
这就是一段上下翻转的摆线。$\varphi=0$ 时位于 $B=(0,y_0)$，$\varphi=\pi$ 时到达
$A=(\pi y_0/2,0)$。此时 $\theta=0$，即射线垂直岸线到达，而不是斜着撞岸。
最后一句是线性浅水波射线模型的极限描述；实际近岸波幅可能不再远小于水深，出现破碎等现象，不能把理想轨迹继续用于这些区域。
\end{xesolution}

\Needspace{16\baselineskip}\begin{origexamplebox}{自由转动直角水管的稳定偏角}\label{ch13:ex:hose}
一根均匀直角软管的内截面积为 $A=2.0\,\mathrm{cm^2}$，软管自身质量为 $m_h=100\,\mathrm g$，
两段管长均为 $L=40\,\mathrm{cm}$，转动时两段形状不变且始终保持直角。进水端在固定接头 $O$ 处可在竖直平面内自由转动；水从第一段流向弯头，再由第二段射出，
体积流量为 $Q=0.30\,\mathrm{L/s}$。取水的密度 $\rho=1.0\times10^3\,\mathrm{kg/m^3}$、
$g=10\,\mathrm{m/s^2}$，忽略黏性和出口之外水流的重量，求稳定时软管相对“第一段竖直向下、第二段水平向右”构形向后转过的角度 $\theta$。

\FluidHoseDiagram
\end{origexamplebox}
\Needspace{6\baselineskip}\begin{xesolution}
\ThAHose

图~\ref{ch13:fig:hose}中 $O$ 是转轴，$E$ 是直角弯头，$P$ 是出水口。两段长度相等且转动中保持垂直。
出口平均速度为
\[
 u=\frac QA=\frac{3.0\times10^{-4}}{2.0\times10^{-4}}
 =1.50\,\mathrm{m/s}.
\]
管内水的质量不能遗漏：
\[
 m_w=\rho A(2L)=0.160\,\mathrm{kg},\qquad
 M=m_h+m_w=0.260\,\mathrm{kg}.
\]
取 $x$ 向右、$y$ 向上。顺时针转过 $\theta$ 后，第一、第二段的单位方向分别为
\[
 \vect e_1=(-\sin\theta,-\cos\theta),\qquad
 \vect e_2=(\cos\theta,-\sin\theta).
\]
软管与其中水的线密度都沿两段均匀分布，整体质心位置为
\[
 \vect r_C=\frac{3L}{4}\vect e_1+\frac L4\vect e_2,\qquad
 x_C=\frac L4(\cos\theta-3\sin\theta).
\]

以软管连同管内的稳态水流为控制体，对 $O$ 点列角动量矩方程。入口就在 $O$，没有角动量矩；出口位置为
$L(\vect e_1+\vect e_2)$，速度为 $u\vect e_2$，故流出的角动量流率为
\[
 \rho Q\,[L(\vect e_1+\vect e_2)]\times(u\vect e_2)
 =\rho QuL\,\vect e_z.
\]
接头反力对 $O$ 无力矩，入口压强力作用线通过 $O$，出口取大气表压为零。因此外力矩只剩总重力矩：
\[
 Mg\frac L4(3\sin\theta-\cos\theta)=\rho QuL.
\]
约去 $L$ 并代入数据，
\[
 3\sin\theta-\cos\theta
 =\frac{4\rho Qu}{Mg}=0.6923.
\]
写成
$\sqrt{10}\sin[\theta-\arctan(1/3)]=0.6923$，取图示的小于 $90^\circ$ 的稳定构形，得到
\[
 \boxed{\theta=31.1^\circ}.
\]
\FluidHoseAnalysisDiagram

这个角度必须大于 $\arctan(1/3)=18.4^\circ$，此后质心位于转轴左侧，重力才提供与流出角动量同向的正力矩；
若误把管内 $0.160\,\mathrm{kg}$ 的水重漏掉，会得到约 $53.1^\circ$，与题给截面积信息的作用不相容。

再作三项核验。第一，等式两边都已除以 $L$：左边 $Mg$ 与右边 $\rho Qu$ 的量纲同为力，
所以无量纲方程没有漏掉长度因子。第二，令
\[
 f(\theta)=3\sin\theta-\cos\theta.
\]
在题图允许的 $0<\theta<90^\circ$ 内，$f'(\theta)=3\cos\theta+\sin\theta>0$，
故 $f$ 单调增加；又有 $f(0)=-1$、$f(90^\circ)=3$，右端 $0.6923$ 落在该区间内，
因此这里恰有一个几何上可接受的解，不会因反三角函数的另一分支产生歧义。

第三，若把流量逐渐降到零，平衡式退化为
$3\sin\theta-\cos\theta=0$，即 $\theta\to\arctan(1/3)$。
这正是直角管连同管内静水仅在重力作用下使质心落到转轴正下方的构形；随着 $Q$ 增大，
出口角动量流率 $\rho Q^2L/A$ 增大，所需重力矩随之增大，因而 $\theta$ 应继续增大。
这个极限与单调性均和求得的 $31.1^\circ$ 一致。
\end{xesolution}

\endgroup

\hypersetup{bookmarksdepth=2}
